\documentclass[twoside, openright, fontsize=10pt, headings=big]{scrbook}

% Load your existing preamble
\input{../../../../preambles/textbook-notes-preamble2.tex}
\usetikzlibrary{cd}
% KOMA header configuration
\usepackage{scrlayer-scrpage}
\clearpairofpagestyles
\automark[chapter]{chapter}
\ohead{\pagemark}
\ihead{\headmark}
\cehead{\MakeUppercase{\TextbookTitleShort}}
\renewcommand{\chapterpagestyle}{empty}
\newcommand{\addunumberedsubsec}[1]{%
    \subsection*{#1}%
    \addcontentsline{toc}{subsection}{#1}%
}

% Define title page fields
\def\YourName{Deepak Jassal}
\def\TextbookTitle{Analytic Number Theory}
\def\TextbookTitleShort{ANT}
\def\TextbookAuthor{Dr. Alia Hamieh}
\def\CourseName{MATH481}
\def\NoteSemester{Winter}
\def\NoteYear{2026}

% Make equations share the exact same counter as theorems
%\makeatletter
%\let\c@equation\c@tcbtheorem
%\renewcommand{\theequation}{\thetcbtheorem}
%\makeatother


\begin{document}
%No page spaces between chapters
\let\cleardoublepage\clearpage
% Title page
% ============================================================
%  titlepage.tex  —  Detailed Textbook Notes Title Page
% ============================================================

\begin{titlepage}
  \centering
  
  \vspace*{1.5cm}
  
  % Top decoration
  {\Large \scshape \textbf{Personal Study Notes} \par}
  
  \vspace{0.3cm}
  \rule{0.3\textwidth}{0.4pt}
  \vspace{1cm}
  
  % Course info
  {\Large \CourseName \par}
  
  \vspace{0.5cm}
  
  % Main title
  {\Large \textsc{Based on} \par}
  {\Huge \bfseries \TextbookTitle \par}
  
  \vspace{0.5cm}
  
  % Textbook author
  {\Large \itshape by \TextbookAuthor \par}
  
  \vspace{1.5cm}
  
  \rule{0.5\textwidth}{0.4pt}
  
  \vspace{1cm}
  
  % Your info
  {\Large \textbf{Notes by:} \YourName \par}
  
  \vspace{0.3cm}
  
  % Date info
  {\large \NoteSemester\ \NoteYear \par}
  
  \vfill
  
  % Footer
  \rule{0.8\textwidth}{0.4pt}
  \vspace{0.2cm}
  
  {\small \textit{These notes are for personal use only.}}
  
\end{titlepage}

\newpage
\thispagestyle{empty}
\mbox{}
\newpage

% ============================================================
%  END OF TITLE PAGE
% ============================================================

\frontmatter
\tableofcontents
\mainmatter

\chapter{Introduction}
Motivating questions in analytic number theory.\\
What is the ``probability'' that a randomly chosen integer is 
\begin{itemize}
    \item prime?
    \item square free?
    \item has an odd number of prime factors?
\end{itemize}
Given a multiplicative function $(f:\N\to\C,\,f(mn)=f(m)f(n),\,\gcd(m,n=1))$, $\varphi(n),\,d(n),\,\sigma(n),$ etc.\\
What is the distribution of such functions on ``randomly'' chosen integers
\begin{itemize}
    \item minimum/maximum
    \item average value
    \item typical value
\end{itemize}
To make some of these questions more rigorous:\\
Let $x\geq 1$ be a real parameter\\
The answer for the questions for $n\in\{1,2,3,\dots,\left\lfloor x\right\rfloor \}$ will depend on our choice of $x$.\\
For example, consider $\dfrac{1}{x}\sum_{n\leq x}f(n)$ where $f(n)$ is a multiplicative function, and study it's behaviour in terms of $x$ for large values of $x$.
\begin{example}{}{}
    The average value $d(n)$
\[
    \lim_{x\to\infty}\frac{1}{x}\sum_{n\leq x}d(n)=\log(x), \quad \frac{1}{x}\sum_{n\leq x}d(n)\sim\log(x) 
\]
\end{example}
Let $\pi(x)$ be the number of prime numbers less than or equal to $x$.
\[
    \pi(x)=\sum_{\substack{p\leq x\\ p\text{ prime}}}1.
\]
The (weak form) of the prime number theorem states that:
\[
    \pi(x)\sim\frac{x}{\log(x)}
\]
for large values of $x$.\\
In other words that is 
\[
    \lim_{x\to\infty}\frac{\pi(x)}{\frac{x}{\log(x)}}=1.
\]
A stronger version of the prime number theorem asserts that 
\[
    \pi(x)=\Li(x)+\Oh\left(e^{-c\sqrt{\log(x)}}\right)
\]
where
\[
    \Li(x)=\int_{2}^{x}\frac{dt}{\log(t)},
\]
is the logarithmic integral function.\\\
Using integration by parts we can see that 
\[
    \Li(x)=\frac{x}{\log(x)}+\Oh\left(\frac{x}{(\log(x))^3}\right).
\]
For a review of arithmetic and multiplicative functions please refer to chatper 1 of Hilderbrand's notes.\\
\setcounter{section}{1}
\chapter{Asymptotic Notation \\\& Average Values of Arithmetic Functions}
\addsec{``Big Oh'' Notation}
Let $f:A\to\C$ be a function, and let $g(x)$ be a non-negative function.\\
We write $f(x)=\Oh(g(x))$ or $f(x)<<g(x)$ is there exists $C>0$ such that
\[
    |f(x)|\leq Cg(x)
\]
for all $x\in A$.\\
In this course we will mostly be interested in teh asymptotic behaviour of a function $f(x)$ for large values of $x$ (i.e., $x\to\infty$). In this case, the ``big oh'' notation can be understood as follows
\[
    f(x)=\Oh(g(x))
\]
whenever there exists $C>0$ and $x_0>0$ such that
\[
    |f(x)|\leq Cg(x)
\]
for all $x\geq x_0$.
\addsec{``Little Oh'' Notation}
We say $f(x)=\oh(g(x))$ if 
\[
    \lim_{x\to\infty}\frac{f(x)}{g(x)}=0
\] provided that $g(x)>0$ for large enough $x$.
In other words, given $\varepsilon>0$, there exists $x_0>1$ such that 
\[
    \left|\frac{f(x)}{g(x)}\right|\leq \varepsilon
\]
for all $x\geq x_0$.
Rearranging this expression we see that 
\[
    |f(x)|\leq \varepsilon g(x)\Leftrightarrow f(x)=\Oh(g(x).)
\]
\begin{example}{}{}
    \begin{enumerate}
        \item $\log(x)=\Oh(x)$ because $\lim_{x\to\infty}\frac{\log(x)}{x}=0$.
        \item $\lim_{x\to\infty}f(x)=0\Leftrightarrow f(x)=\Oh(1)$. 
    \end{enumerate}
\end{example}
\addsec{Asymptotic Equivalence}
We write
\[
    f(x)\sim g(x)
\]
if 
\[
    \lim_{x\to\infty}\frac{f(x)}{g(x)}=1.
\]
\addsec{Order of Magnitude}
We write 
\[
    f(x)\asymp g(x)
\]
if
\[
    f(x)<<g(x),\text{ and} g(x)<<f(x).
\]
In other words, there exists $c_1,c_2\geq 0$ such that
\[
    c_1g(x)\leq f(x)\leq c_2g(x)
\]
\begin{remark}
    $f(x)\sim g(x) \Rightarrow f(x)\asymp g(x)$ but the backwards implication is not always true.
\end{remark}
\begin{example}{}{}
    $x^2\asymp x^2,$ but $\lim_{x\to\infty}\frac{2x^2}{x^2}=2\neq1.$
\end{example}
\addsec{Asymptotic Formula}
\[
    f(x)=\underbrace{g(x)}_{\text{main term}}+\underbrace{\Oh(\mathcal{R}(x))}_{\text{error term}}
\]
\[
    |f(x)-g(x)|=\Oh(\mathcal{R}(x))\Leftrightarrow \exists C>0 \text{ such that} |f(x)-g(x)|\leq C\Oh(\mathcal{R}(x)).
\]
An asymptotic formula is only meaningful if 
\[
    \mathcal{R}(x)=\oh(g(x))
\]
i.e., the error term is much smaller than the main term.
\begin{example}{}{}
    $\cos(x)=1+\Oh(x^2)$ for $|x|<1$.
\end{example}
\begin{proof}
    By the mean value theorem
    \[
    \left|\frac{\cos(x)-\cos(0)}{x-0}\right|=|\sin(c)|
    \]
    for some $0\leq c<|x|$. From here we can see that
    \[
    \left|\frac{1-\cos(x)}{|x|}\right|=|\sin(c)|\leq \max_{0\leq y\leq |x|}\{\sin(y)\}\leq |x|.
    \]
    Hence, 
    \begin{align*}
        1-\cos(x)&\leq x^2\\
        &=\Oh(x^2)\\
        \cos(x)&=1+\Oh(x^2).\qedhere 
    \end{align*}
\end{proof}
\addsec{Dependance on Parameters}
Suppose that $f$ and $g$ depend on a parameter $\lambda$. Then, the notation 
\[
    f(x)=\Oh_{\lambda}(g(x))
\]
\[
    f(x)\underset{\lambda}{<<}(g(x))
\]
means that the implicit constant depends on $\lambda$.
\begin{example}{}{}
\begin{enumerate}
    \item $\log(x)=\Oh_\varepsilon(x^\varepsilon)$ for any $\varepsilon>0$.
    \item Given $0<\alpha<1$ then for any $A>0$ and $\varepsilon>0$ we have
    \[
        x^{-\varepsilon}\underset{\alpha,\varepsilon}{<<}e^{-c(\log(x))^\alpha}\underset{\alpha,A}{<<} (\log(x))^{-A}.
    \]
\end{enumerate}
\end{example}
\addsec{Next Time}
Next time we will consdier 
\[
    \frac{1}{x}\sum_{n\leq x}\mu(n),
\]
and then move on to developing the summation formulas that allow us to compute average values ($\frac{1}{x}\sum{n\leq x}f(n)$) of arithmetic functions.

Many important problems in analytic number theory reduce to evaluating (partial) sums of the form
\[
    \sum_{n\leq x}a_n
\]
where $a_n$ is an arithmetic function.
\begin{example}{}{}
    \begin{enumerate}
        \item Let $\Psi(x)=\sum_{n\leq x}\Lambda(n)$ where $\Lambda(n)$ is the Von Mangoldt function. The prime number theorem (PNT) is equivalent to $\Psi(x)\sim x$.\\ Under the Reimann hypothesis we have
        \[
            \Psi(x)=x+\Oh\left(x^{\frac{1}{2}+\varepsilon}\right)
        \]
        for any $\varepsilon>0$.
        \item The PNT is equivalent to 
        \[
            \sum_{n\leq x}\mu(n)=\oh(x).
        \]
        i.e.,
        \[
            \lim_{x\to\infty}\frac{1}{x}\sum_{n\leq x}\mu(n)=0.
        \]
        Where $\mu(n)$ is the M\"{o}bius function.\\
        Some interesting facts about $\sum_{n\leq x}\mu(n)$. Merten's conjectured that $|\sum_{n\leq x}\mu(n)|\leq\sqrt{x}$. In 1985 Odlyzko proved that $|\sum_{n\leq x}\mu(n)|>1.06\sqrt{x}$ for infinitely many $x$.\\
        $\sum_{n\leq x}\mu(n)$ is not a straighforward or easy sum to work with and get significant results on it's asymptotic growth.\\
        Instead of considering this sum we will consider the weighted sum
        \[
        \sum_{n\leq x}\frac{\mu(n)}{n}.
        \]
        This is a common trick in analytic number theory (ANT). If $\sum_{n\leq x}f(n)$ is a difficult sum to work with, instead work with $\sum_{n\leq x}\frac{\mu(n)}{\omega(n)}$ where $\omega(n)$ is a ``nice'' weight function.
        \item Special Case of the Hardy Littlewood Conjecture
        \[
            \sum_{n\leq x}\Lambda(n)\Lambda(n+2)\sim Cx
        \]
        Where $C$ is given by \[C=\prod_{\substack{p\geq 3\\p\text{ prime}}}\left(1-\frac{1}{(p-1)^2}\right).\]
        This is a special case of the twin prime conjecture.
    \end{enumerate}
\end{example}
Let's consider some partial suns that bac be dealt with without the use of summation formulas.
\begin{example}{}{}
    \begin{enumerate}
        \item $\sum_{n\leq x}n=\left\lfloor x\right\rfloor$, where $\left\lfloor x\right\rfloor=x-\{x\}=x+\Oh(x)$ and $\{x\}$ is the fractional part of $x$.
        \item 
        \begin{align*}
            \sum_{n\leq x}d(n)&=\sum_{n\leq x}\sum_{d\mid n}1\\
            &=\sum_{d\leq x}\sum_{n\leq x}1\\
            &=\sum_{d\leq x}\left\lfloor\frac{x}{d}\right\rfloor\\
            &=\sum_{d\leq x}\left(\frac{x}{d}+\Oh(x)\right)\\
            &=x\sum_{d\leq x}\frac{1}{d}+\Oh(x)\\
            &=x(\log x +\Oh(x))+\Oh(x)\\
            \sum_{n\leq x}d(n)&=x\log x +\Oh(x).
        \end{align*}
        From here it is clear to see that
        \[
            \frac{1}{x}\sum_{n\leq x}d(n)=\log x + \Oh(1).
        \]
        \item We prove that
        \[
            \left|\sum_{n\leq x}\frac{\mu(n)}{n}\right|\leq 1.
        \]
        We have
        \begin{align*}
            \left|\sum_{n\leq x}\frac{\mu(n)}{n}\right|&=\frac{1}{x}\left|\sum_{n\leq x}\mu(n)\frac{x}{n}\right|\\
            &=\frac{1}{x}\left|\sum_{n\leq x}\mu(n)\left(\left\lfloor \frac{x}{n}\right\rfloor + \left\{\frac{x}{n}\right\}\right)\right|\\
            &\leq\frac{1}{x}\left|\sum_{n\leq x}\mu(n)\left\lfloor \frac{x}{n}\right\rfloor\right| + \frac{1}{x}\left|\sum_{n\leq x}\mu(n)\left\{ \frac{x}{n}\right\}\right|.
        \end{align*}
        We will first deal with the sum on the right
        \begin{align*}
            \frac{1}{x}\left|\sum_{n\leq x}\mu(n)\left\{ \frac{x}{n}\right\}\right|&\leq\frac{1}{x}\sum_{n\leq x}\left|\mu(n)\right|\left\{ \frac{x}{n}\right\}\\
            &\leq\sum_{n\leq x}\left\{\frac{x}{n}\right\}\\
            &=\sum_{n\leq \lfloor x\rfloor}\left\{\frac{x}{n}\right\}\\
            &=\sum_{n\leq \lfloor x\rfloor -1}\left\{\frac{x}{n} \right\}+\left\{\frac{x}{\lfloor x\rfloor} \right\}\\
            &\leq\floor{x}-1+\left\{\frac{\floor{x}+\{x\}}{\floor{x}}\right\}\\
            &\leq \floor{x}-1+\left\{1+\frac{\{x\}}{\floor{x}} \right\}\\
            &\leq \floor{x}-1+\frac{\fract{x}}{\floor{x}}\\
            &\leq \floor{x}-1+\fract{x}=x-1.
        \end{align*}
        Now for the left hand sum
        \begin{align*}
            \sum_{n\leq x}\mu(n)\floor{\frac{x}{n}}&=\sum_{n\leq x}\mu(n)\sum_{m\leq \frac{x}{n}\Leftrightarrow l=mn\leq x}1\\
            &\sum_{l\leq x}\sum_{n\mid l}\mu(n)\\
            &=1
        \end{align*}
        From MATH480 we know that 
        \[
            \sum_{n\mid l}\mu(n)=
            \begin{cases}
                1& \text{if } l=1\\
                0& \text{otherwise}.
            \end{cases}
        \]
        Combining these two results we arrive at
        \begin{align*}
            \left|\sum_{n\leq x}\frac{\mu(n)}{n}\right|&=\frac{1}{x}\left|\sum_{n\leq x}\mu(n)\frac{x}{n}\right|\\
            &\leq \frac{1}{x}(1+x-1)\\
            =1.
        \end{align*}
    \end{enumerate}
\end{example}
\addsec{Summation of Smooth Functions (Euler's Summation Formula)}
\begin{theorem}{Euler's Summation Formula}{eulersum}
    Let $o<y\leq x$ be real numbers and $f\in\mathcal{C}[y,x]$. Then,
    \[
        \sum_{y<n\leq x}f(n)=\int_{y}^{x}f(t)\,dt+\int_{y}^{x}\fract{t}f'(t)dt-\fract{x}f(x)+\fract{y}f(y).
    \]
\end{theorem}
\begin{corollary}{}{}
    Let $x\geq 1$ be a real number, $f\in\mathcal{C}[1,x]$. Then \ref{thm:eulersum} gives
    \[
        \sum_{n\leq x}f(n)=\int_{1}^{x}f(t)\,dt+\int_{1}^{x}\fract{t}f'(t)dt-\fract{x}f(x)+f(1)
    \]
\end{corollary}
\begin{proof}
    Apply theorem \ref{thm:eulersum} with $y=1$.
\end{proof}
Before proving theorem \ref{thm:eulersum} we establish the following application which was stated earlier.
\begin{theorem}{Partial Sums of the Harmonic Series}{harmonic}
    \[
        \sum_{n\leq x}\frac{1}{n}=\log x +\gamma +\Oh\left(\frac{1}{x}\right)
    \]
    where $\gamma$ is the Euler-Mascheroni constant defined asymptotic
    \[
        \gamma=\lim_{x\to\infty}\left(\sum_{n\leq x}\frac{1}{n}-\log x\right)=0.5772\dots
    \]
\end{theorem}
\begin{proof}
    Applying corollary 2.2 with $f(n)=\frac{1}{n}$ we get the following
    \begin{align*}
        \sum_{n\leq x}\frac{1}{n}&=\int_{1}^{x}\frac{1}{t}\,dt+\int_{1}^{x}\fract{t}\left(\frac{-1}{t^2}\right)\,dt-\fract{x}\frac{1}{x}+1\\
        &=\log x +1 -\int_{1}^{x}\frac{\fract{t}}{t^2}\,dt+\Oh\left(\frac{1}{x}\right).
    \end{align*}
    Set $I(x)=\int_{1}^{x}\frac{\fract{t}}{t^2}\,dt$. Then,
    \begin{align*}
        I(x)&=\int_{1}^{\infty}\frac{\fract{t}}{t^2}\,dt-\int_{x}^{\infty}\frac{\fract{t}}{t^2}\,dt\\
        &=I-\int_{x}^{\infty}\frac{\fract{t}}{t^2}\,dt.
    \end{align*}    
    Observe that
    \[
        \int_{x}^{\infty}\frac{\fract{t}}{t^2}\,dt\leq \int_{x}^{\infty}\frac{1}{t^2}\,dt=\left[\frac{-1}{t}\right]^\infty_x=\frac{1}{x}.
    \]
    Thus,
    \[
        I(x)=I+\Oh\left(\frac{1}{x}\right).
    \]
    We now have
    \[
        \sum_{n\leq x}\frac{1}{n}=\log x -I + 1+\Oh\left(\frac{1}{x}\right).    
    \]
    Rearranging this we see that
    \[
    \lim_{x\to\infty}\left(\sum_{n\leq x}\frac{1}{n}-\log x\right)=\gamma=\lim_{x\to\infty}\left(1-I+\Oh\left(\frac{1}{x}\right)\right).
    \]
    Hence,
    \[
        \sum_{n\leq x}\frac{1}{n}=\log x +\gamma +\Oh\left(\frac{1}{x}\right). \qedhere
    \]
\end{proof}
\begin{proof}[Proof of Theorem \ref{thm:eulersum}]
    Set $k=\floor{x}$ and $m=\floor{y}$. Then we rewrite the sum asymptotic
    \[
        \sum_{y<n\leq x}f(n)=\sum_{m+1\leq n\leq k}f(n).
    \]
    Consider $n$ such that $n-1,n\in[y,x]$. Then,
    \begin{align*}
    \int_{n-1}^{n}\floor{t}f(t)\,dt&=\int_{n-1}^{n}(n-1)f'(t)\,dt\\
    &=(n-1)(f(n)-f(n-1))\\
    &=nf(n)-(n-1)f(n-1)-f(n).\qedhere 
    \end{align*}
\end{proof}

\textit{Recall.} Eulers summation formula which states that for $f\in\mathcal{C}[y,x]$ we have
\[
    \sum_{y\leq n\leq x}f(n)=\int_{y}^{x}f(t)\,dt+\int_{x}^{y}\fract{t}f'(t)\,dt-\fract{t}f(x)+\fract{y}f(y).
\]
\begin{proof}[Continuation of Proof for Theorem 2.1]
    Set $k=\floor{x}$ and $m=\floor{y}$. We observe that
    \[
        \sum_{y\leq n\leq x}f(n)=\sum_{m+1\leq n\leq k}f(n).
    \]
    We also have for $n-1,n\in[y,x]$
    \[
        \int_{n-1}^{n}\floor{t}f'(t)\,dt=nf(n)-(n-1)f(n-1)-f(n).
    \]
    Hence,
    \begin{align*}
        \int_{m+1}^{k}\floor{t}f'(t)\,dt=&\int_{m+1}^{m+2}\floor{t}f(t)\,dt+\int_{m+2}^{m+3}\floor{t}f(t)\,dt+\cdots+\int_{k-1}^{k}\floor{t}f(t)\,dt\\
        =&(m+2)f(m+2)-(m+1)f(m+1)-f(m+2)\\
        &+(m+3)f(m+3)-(m+2)f(m+1)-f(m+3)\\
        &+\;\,\vdots\\
        &+kf(k)-(k-1)f(k-1)-f(k)\\
        =&-mf(m+1)+kf(k)-\sum_{m-1\leq n\leq k}f(n).
    \end{align*}
    It follows that
    \begin{align*}
        \sum_{m-1\leq n\leq k}f(n)&=kf(k)-mf(m+1)-\int_{m+1}^{k}\floor{k}f'(t)\,dt\\
        &=kf(k)-mf(m+1)-\int_{y}^{x}\floor{t}f'(t)\,dt+\int_{y}^{m+1}\floor{t}f'(t)\,dt+\int_{k}^{x}\floor{t}f'(t)\,dt\\
        &=kf(k)-mf(m+1)-\int_{y}^{x}\floor{t}f'(t)\,dt+m\int_{y}^{m+1}f'(t)\,dt+k\int_{k}^{x}f'(t)\,dt\\
        &=kf(k)-mf(m+1)-\int_{y}^{x}\floor{t}f'(t)\,dt+m(f(m+1)-f(y))+k(f(x)-f(k))\\
        &=kf(x)-mf(y)-\int_{y}^{x}tf'(t)\,dt+\int_{y}^{x}\fract{t}f'(t)\,dt\\
        &=kf(x)-mf(y)-[tf(t)]_y^x+\int_{x}^{y}f(t)\,dt+\int_{y}^{x}\fract{t}f'(t)\,dt\\
        &=xf(x)-\fract{x}f(x)-yf(y)+\fract{y}f(y)-xf(x)+yf(y)+\int_{y}^{x}f(t)\,dt+\int_{y}^{x}\fract{t}f(t)\,dt.
    \end{align*}
    Cancelling the terms we arrive at 
    \[
                \sum_{y\leq n\leq x}f(n)=\int_{y}^{x}f(t)\,dt+\int_{x}^{y}\fract{t}f'(t)\,dt-\fract{t}f(x)+\fract{y}f(y).\qedhere
    \]
\end{proof}
\begin{theorem}{Sum of Logartihm's}{}
    Let $N\in\N$ then we have
    \[
    \sum_{n\leq N}\log n=Nlog N-N+\frac{1}{2}\log N + c+\Oh\left(\frac{1}{N}\right).
    \]
\end{theorem}

Observe that the above formula can be written as
\[
    \log N!=Nlog N-N+\frac{1}{2}\log N + c+\Oh\left(\frac{1}{N}\right),
\]
exponentiating we obtain
\begin{equation}\label{eqn:stirling}
    N!=N^Ne^{-N}\sqrt{N}e^ce^{\Oh\left(\frac{1}{N}\right)}.
\end{equation}
What can be said about $e^{\Oh\left(\frac{1}{N}\right)}$?\\
Let $g(N)$ be a function such that $g(N)=\Oh\left(\frac{1}{N}\right)$. i.e., $|g(N)|\leq C\frac{1}{N}$ for some $C>0$.
\begin{align*}
    \left|e^{g(N)}\right|&=\left|a+g(N)+\frac{g(N)^2}{2!}++\frac{g(N)^3}{3!}+\cdots\right|\\
    &\leq1+\left|g(N)\right|+\left|\frac{g(N)^2}{2!}\right|+\cdots\\
    &\leq1+\frac{C}{N}+\frac{C^2}{N}\\
    &=1+\frac{C}{N}\left(1+\frac{1}{2!}\frac{C}{N}+\frac{1}{3!}\frac{C^2}{N^2}+\cdots\right)\\
    &=1+\Oh\left(\frac{1}{N}\right).
\end{align*}
It follows that $e^{\Oh\left(\frac{1}{N}\right)}=1+\Oh(\frac{1}{N})$. Furthermore,
\[
    e^x=1+\Oh\left(x\right),\quad \text{if } |x|<1.
\]
Going back to equation \ref{eqn:stirling} we have
\[
    N!=N^Ne^{-N}\sqrt{N}e^c\left(1+\Oh\left(\frac{1}{N}\right)\right).
\]
Hence, we arrive at Stirling's formula.
\begin{corollary}{Stirling's formula}{}
    \[
        N!=C_1N^Ne^{-N}\sqrt{N}\left(1+\Oh\left(\frac{1}{N}\right)\right).
    \]   
\end{corollary}
\begin{proof}
    By Euler's summation formula we have
    \begin{align*}
        \sum_{n\leq N}\log n&=\int_{1}^{N}\log t\,dt+\int_{1}^{t}\fract{t}\frac{1}{t}\,dt-\fract{N}\log N\\
        &=N\log N-N+\int_{1}^{N}\frac{\fract{t}}{t}\,dt.
    \end{align*}
    Now let's consider $\int_{1}^{N}\frac{\fract{t}}{t}\,dt$.
    \begin{align*}
        \int_{1}^{N}\frac{\fract{t}}{t}\,dt&=\int_{1}^{N}\frac{\fract{t}+\frac{1}{2}-\frac{1}{2}}{t}\,dt\\
        &=\frac{1}{2}\log N+\int_{1}^{N}\frac{\fract{t}-\frac{1}{2}}{t}\,dt.
    \end{align*}
    Out problem is now reduced to showing that $\int_{1}^{N}\frac{\fract{t}-\frac{1}{2}}{t}\,dt=c+\Oh\left(\frac{1}{N}\right)$, for some explicitly given constant $c$.\\
    Let $\rho(t)=\frac{t}-\frac{1}{2}$, then we have
    \[
        \int_{1}^{N}\frac{\rho(t)}{t}\,dt
    \]
    let $u=\frac{1}{t}$, $dv=\rho(t)\,dt$, $du=\frac{1}{t^2}\,dt$, $v=\int_{1}^{t}\rho(x)\,dx=R(t)$. Then the above integral becomes
    \[
        \left[\frac{R(t)}{t}\right]_1^N+\int_{1}^{N}\frac{1}{t^2}\,dt.
    \]
    Notice a property of $R(t)$.
    \begin{align*}
        R(N)=\int_{1}^{N}\rho(x)\,dx&=\int_{1}^{2}\rho(x)\,dx+\int_{2}^{3}\rho(x)\,dx+\cdots+\int_{N-1}^{N}\rho(x)\,dx\\
        &=\int_{1}^{2}\fract{x}-\frac{1}{2}\,dx+\int_{2}^{3}\fract{x}-\frac{1}{2}\,dx+\cdots+\int_{N-1}^{N}\fract{x}-\frac{1}{2}\,dx\\
        &=\int_{0}^{1}x-\frac{1}{2}\,dx+\int_{0}^{1}x-\frac{1}{2}\,dx+\cdots+\int_{0}^{1}x-\frac{1}{2}\,dx\\
        &=\left[\frac{x^2}{2}-\frac{x}{2}\right]_0^1+\left[\frac{x^2}{2}-\frac{x}{2}\right]_0^1+\cdots+\left[\frac{x^2}{2}-\frac{x}{2}\right]_0^1\\
        &=0+0+\cdots+0\\
        &=0.
    \end{align*}
    Hence, $R(N)=0$ for all $N\in\N$.\\
    Observe that
    \begin{align*}
        \left|R(t)\right|&=\left|\int_{1}^{t}\rho(x)\,dx\right|\\
        &=\left|\int_{\floor{t}}^{t}\rho(x)\,dx\right|\\
        &\leq\int_{\floor{t}}^{t}\left|\rho(x)\right|\,dx\\
        &\leq \frac{1}{2}\left(t-\floor{t}\right)=\frac{1}{2}\fract{t}\leq\frac{1}{2}.
    \end{align*}
    Hence, $\int_{1}^{\infty}\frac{R(t)}{t^2}\,dt=c$ for some constant $c$, and 
    \[
        \int_{N}^{\infty}\frac{R(t)}{t^2}\,dt\leq \int_{N}^{\infty}\frac{\frac{1}{2}}{t^2}\,dt\ll\frac{1}{N}.
    \]
    Thus we have
    \[
        \sum_{n\leq N}\log n =N\log N-N+\frac{1}{2}\log N+c+\Oh\left(\frac{1}{N}\right).\qedhere
    \]
\end{proof}
\begin{corollary}{Sum of Logarithm's}{}
    Given $x\in\R$ and $x\geq1$ we have
\[
    \sum_{n\leq N}\log n=x\log x-x+\Oh(\log x).
\]
\end{corollary}
\begin{proof}
    The result can be derived directly from Euler's summation formula, we will instead use another approach.\\
    Applying theorem 2.4 with $N=\floor{x}$
    \begin{align*}
        \sum_{n\leq x}\log n&=\sum_{n\leq N}\log n\\
        &=\floor{x}\log\floor{x}-\floor{x}+c+\frac{1}{2}\log\floor{x}+\Oh\left(\frac{1}{\floor{x}}\right)\\
        &=(x-\fract{x})\log(x-\fract{x})-(x-\fract{x})+c+\frac{1}{2}\log(x-\fract{x})+\Oh\left(\frac{1}{(x-\fract{x})}\right)
    \end{align*}
    Given $|t|<1$
    \begin{align*}
        \log t&=-\sum_{n=1}^{\infty}\frac{t^n}{n}=1-t-\frac{t^2}{2}-\cdots\\
        &=-1-t\left(1+\frac{t}{2}+\frac{t^2}{3}+\cdots\right)=-1+\Oh(t).
    \end{align*}
    Since $|t|<1$ we further have $\log t=\Oh(1)$.\\
    Hence,
    \[
        \log(x-\fract{x})=\log x+\Oh(1).\qedhere
    \]
\end{proof}

Another application of Euler's summation formula.\\
The Reimann zeta function is defined as 
\[
    \zeta(s)=\sum_{n=1}^{\infty}\frac{1}{n^s},\quad s\in\C,\; \Re(s)>1.
\]
This is defined as an analytic function in the complex variable $s$ on the half plane $\Re(s)>1$. $\zeta(s)$ converges absolutely for $\Re(s)>1$.
\begin{example}{}{}
    \[
        \zeta(2)=\sum_{n=1}^{\infty}\frac{1}{n^2}=\frac{\pi^2}{6}
    \]
\end{example}
\begin{example}{}{}
    \[
        \zeta(2k)\in\Q^{2k},\quad k\in\N.
    \]
\end{example}
\begin{theorem}{Integral Form of the Reimann Zeta Function}{}
    \[
        \zeta(s)=\frac{s}{s-1}-s\int_{1}^{\infty}\frac{\fract{x}}{x^{s+1}}\,dx,\quad \Re(s)>0.
    \]
\end{theorem}
\begin{proof}[Proof of Theorem 2.7]
    Consider for $x\geq1$
    \[
        S(x)=\sum_{n\leq x}\frac{1}{n^s},
    \]
    observe that 
    \[
        \lim_{x\to\infty}S(x)=\zeta(s).
    \]
    By Euler's summation formula we have
    \begin{align*}
        S(x)&=\sum_{n\leq x}\frac{1}{n^s}\\
        &=\int_{1}^{x}\frac{1}{t^s}\,dt-s\int_{1}^{x}\fract{t}t^{-s-1}\,dt-\fract{x}f(x)+f(1)\\
        &=\frac{1}{s-1}-\frac{x^{-s+1}}{s-1}-s\int_{1}^{\infty}\fract{t}t^{-s-1}\,dt+\int_{x}^{\infty}+\int_{x}^{\infty}\fract{t}t^{-s-1}\,dt-\fract{x}x^{-s}+1.
    \end{align*}
    Note that
    \begin{align*}
        \left|\int_{x}^{\infty}\fract{t}t^{-s-1}\,dt\right|&\leq\int_{x}^{\infty}\fract{t}|t^{-s-1}|\,dt\\
        &=\int_{x}^{\infty}\fract{t}t^{\Re(s)-1}\,dt\\
        &\leq\int_{x}^{\infty}t^{\Re(s)-1}\,dt\\
        &=\left(\frac{t^{-\Re(s)}}{-\Re(s)}\right)^\infty_x=\frac{x^{-\Re(s)}}{-\Re(s)}.
    \end{align*}
    Hence, 
    \[
        S(x)=\frac{1}{s-1}-\frac{x^{-s-1}}{s-1}-s\left(\int_{1}^{\infty}\fract{t}t^{-s-1}\,dt+\Oh\left(\frac{x^{-\Re(s)}}{-\Re(s)}\right)\right).
    \]
    As $x\to\infty$ we get $\zeta(s)$
    \begin{align*}
        \zeta(s)&=\frac{1}{s-1}-s\int_{1}^{\infty}\fract{t}t^{-s-1}\,dt+1\\
        &=\frac{s}{s-1}-s\int_{1}^{\infty}\fract{t}t^{-s-1}\,dt.\qedhere
    \end{align*}
\end{proof}
Recall PNT
    \[
        \pi(x)=
    \]

\begin{lemma}{Logarithmic Integral}{}
    \[
        \mathrm{Li}_k(x)=\int_{2}^{x}\frac{1}{(\log t)^k}\,dt\ll_k\frac{x}{(\log x)^k}.
    \]
\end{lemma}
\begin{proof}
    \[
        \int_{2}^{x}\frac{1}{(\log x)^k}\,dt
    \]
    we have
    \[
        \log t\geq\log2,\quad \frac{1}{\log t}\leq\frac{1}{\log2},\quad \frac{1}{(\log t)^k}<\frac{x-2}{(\log2)^k}\ll x-2\ll x.
    \]
    \[
        \underbrace{\int_{2}^{\sqrt{x}}\frac{dt}{(\log t)^k}}_{I_1}+\underbrace{\int_{\sqrt{x}}^{x}\frac{dt}{(\log t)^k}}_{I_2}.
    \]
    \[
        I_1\leq\int_{2}^{\sqrt{x}}\frac{dt}{(\log2)^k}=\frac{\sqrt{x}-2}{(\log2)^k}\ll_k\sqrt{x}\ll\frac{x}{(\log x)^k},
    \]
    for $\sqrt{x}< t< x$
    \[
        \log t\geq \log\sqrt{2}=\frac{1}{2}\log x
    \]
    \[
        \frac{1}{(\log t)^k}\leq\frac{2^k}{(\log x)^k}.
    \]
    \[
        I_2\leq \int_{\sqrt{x}}^{x}\frac{2^k}{(\log t)^k}\,dt=\frac{2^k(x-\sqrt{x})}{(\log x)^k}\ll_k\frac{x}{(\log x)^k}.
    \]
    Hence, 
    \[
        \int_{2}^{x}\frac{2^k}{(\log t)^k}\,dt=\Oh\left(\frac{x}{(\log x)^k}\right),\quad x\geq 4.
    \]
    For $2\leq x\leq 4$ we have
    \[
        \int_{2}^{x}\frac{1}{(\log t)^k}\,dt\leq\int_{2}^{x}\frac{1}{(\log 2)^k}\,dt=\frac{x-2}{(\log 2)^k}\leq\frac{x}{(\log 2)^k}\leq\left(\frac{\log x}{\log2}\right)\frac{x}{(\log x)^k}\ll\frac{x}{(\log x)^k}.\qedhere
    \]
\end{proof}
\begin{theorem}{Summation of Logarithmic Integral}{}
    \[
        \mathrm{Li(x)}=\int_{2}^{x}\frac{dt}{\log t}=\frac{x}{\log x}\left(\sum_{i=1}^{k-1}\frac{i!}{(\log x)^i}+\Oh\left(\frac{x}{(\log x)^k}\right)\right).
    \]
\end{theorem}
\begin{proof}[Proof of Theorem 2.9]
    We have 
    \[
        \int_{2}^{x}\frac{dt}{\log t}
    \]
    \[
        u=\frac{1}{\log t},\quad dv=dt,\quad du=\frac{-1}{(\log t)^2},\quad v=t
    \]
    \begin{align*}
        \int_{2}^{x}\frac{dt}{\log t}&=\left[\frac{t}{\log t}\right]^x_2+\int_{2}^{x}\frac{dt}{t(\log t)^2}\\
        &=\frac{x}{\log x}-\frac{2}{\log2}+\frac{x}{(\log x)^2}-\frac{2}{(\log 2)^2}+2\int_{2}^{x}\frac{dt}{(\log t)^3}\\
        &=\frac{x}{\log x}+\frac{x}{(\log x)^2}+\frac{2x}{(\log x)^3}+\cdots+k\int_{2}^{x}\frac{dt}{(\log t)^{k+1}}+\Oh_k(1)\\
        &=\frac{x}{\log x}+\frac{x}{(\log x)^2}+\frac{2x}{(\log x)^3}+\cdots+(k-1)!\frac{x}{(\log x)^k}+\Oh_k\left(\frac{x}{(\log x)^{k+1}}\right)\\
        &=\frac{x}{\log x}\left(\sum_{i=1}^{k-1}\frac{i!}{(\log x)^i}+\Oh\left(\frac{1}{(\log x)^k}\right)\right).\qedhere
    \end{align*}
\end{proof}
We can use Euler's summation formula for sum continuous functions. Abel's Summation Formula allows for us to sum arithmetic functions.
\begin{theorem}{Abel's Summation Formula}{}
    Let $0<y\leq x$, $f\in\mathcal{C}[y,x]$ and $a(n)$ an arithmetic function, then
    \[
        \sum_{y<n\leq x}a(n)f(n)=A(n)f(x)-A(y)f(y)-\int_{y}^{x}A(t)f'(t)\,dt,
    \]
    where $A(n)$ is the summatory fuction for $a(n)$ $\left(A(n)=\sum_{n\leq x}a(n)\right)$.
\end{theorem}
\begin{corollary}{}{}
    Let $x\leq 1$, $f\in\mathcal{C}[1,x]$ and $a(n)$ be an atrithmetic function. Then,
    \[
        \sum_{n\leq x}a(n)f(n)=A(n)f(x)-\int_{y}^{x}A(t)f'(t)\,dt.
    \]
\end{corollary}
\begin{proof}
    Apply \textit{theorem 2.10}.
\end{proof}
\begin{example}{}{}
    Deduce Euler's summation formula from Abel's summation formula. (Use $a(n)=1$, $A(t)=\floor{t}$.)
\end{example}

\begin{proof}[Proof of Theorem 2.10]
    Consider 
    \[
        I=\int_{y}^{x}A(t)f't\,dt=\int_{y}^{x}\sum_{x\leq t}a(n)f't\,dt.
    \]
    Let 
    \[
        \chi(x,n)=\begin{cases}
            1&\text{if } n\leq t\\
            0&\text{otherwise}
        \end{cases}.
    \]
    Using $\chi$ we Obtain
    \begin{align*}
        I&=\int_{y}^{x}\sum_{n\leq x}a(n)\chi(n,t)f'(t)\,dt\\
        &=\sum_{n\leq x}a(n)\int_{y}^{x}\chi(n,t)f'(t)\,dt\\
        &=\sum_{n\leq x}a(n)\int_{\max{n,y}}^{x}f'(t)\,dt\\
        &=\sum_{n\leq x}a(n)\left(f(x)-f(\max{n,y})\right)\\
        &=\sum_{n\leq x}a(n)f(x)-\left(\sum_{n\leq y}a(n)f(y)+\sum_{y<n\leq x}a(n)f(n)\right)\\
        \int_{y}^{x}A(t)f'(t)\,dt&=A(x)f(x)-A(y)f(y)-\sum_{y<n\leq x}a(n)f(n)\\
        \sum_{y<n\leq x}a(n)f(n)&=A(n)f(x)-A(y)f(y)-\int_{y}^{x}A(t)f'(t)\,dt.\qedhere
    \end{align*}
\end{proof}
\begin{theorem}{Kronecker's Lemma}{}
    Let $f:\N\to\C$ be an arithmetic function, let $s\in\C$ with $\Re(s)>0$. Suppose that
    \[
        F(s)=\sum_{n=1}^{\infty}\frac{f(n)}{n^s}
    \]
    converges. Then, 
    \[
        \frac{1}{x^s}\sum_{n\leq x}f(s)\to0\text{ as } x\to\infty.
    \]
    In particular if $\sum_{n=1}^{\infty}\frac{f(s)}{n}$ converges then we have $\frac{1}{x}\sum_{n\leq x}f(s)\to 0$ as $x\to\infty$.\\
    i.e., $\sum_{n\leq x}=\oh(x).$
\end{theorem}
\begin{proof}[Proof of Theorem 2.12]
    Let 
    \[
        S_f(x)=S(x)=\sum_{n\leq x}f(n), 
    \]
    and 
    \[
        T(x)=\sum_{n\leq x}\frac{f(n)}{n^s}.
    \]
    We know that $\lim_{x\to\infty}T(x)=T(s)$. We want to prove that 
    \[
        \lim_{x\to\infty}\frac{S(x)}{x^s}=0.
    \]
    \begin{align*}
        S(x)=\sum_{n\leq x}\frac{f(n)}{n^s}n^s&=\left(\sum_{n\leq x}\frac{f(n)}{n^s}\right)x^s-s\int_{1}^{x}\left(\sum_{n\leq t}\frac{f(n)}{n^s}\right)t^{s-1}\,dt\\
        &=T(x)x^s-s\int_{0}^{x}T(t)t^{s-1}\,dt\\
        &=T(x)\int_{0}^{x}st^{s-1}\,dt-s\int_{0}^{x}T(t)t^{s-1}\,dt\\
        &=\int_{0}^{x}st^{s-1}(T(x)-T(t))\,dt
    \end{align*}
    Hence, since $\lim_{x\to\infty}T(x)=T$ we have
    \[
        |S(x)|\leq\int_{0}^{x}|s|t^{\Re(s)-1}(|T(x)-T|+|T(t)-T|)\,dt.
    \]
    Given $\varepsilon>0$, there exists $x_0>1$ such that $|T(t)-T|<\varepsilon$ whenever $x>x_0$. Given $\varepsilon>0$ we have
    \begin{align*}
        |S(x)|\leq&\int_{0}^{x_0}|s|t^{\Re(s)-1}(|T(x)-T|+|T(t)-T|)\,dt\\
        &+\int_{x_0}^{x}|s|t^{\Re(s)-1}(|T(x)-T|+|T(t)-T|)\,dt.
    \end{align*}
    So for $0<t<x_0$, we have
    \begin{align*}
        |T(x)-T|+|T(t)-T|&\leq\varepsilon +|T(t)|+|T|\\
        &\leq\varepsilon+|T|+\sum_{n\leq t}\frac{|f(n)|}{n^{\Re(s)}}
        &\leq\varepsilon+|T|+\sum_{n\leq x_0}\frac{|f(n)|}{n^{\Re(s)}}\\
        &=M(\varepsilon).
    \end{align*}
    For $x_0<t<x$
    \[
        |T(x)-T|+|T(t)-T|\leq2\varepsilon.
    \]
    Putting the two together we get
    \begin{align*}
        |S(x)|&\leq\int_{0}^{x_0}|s|t^{\Re(s)-1}M(\varepsilon)\,dt+\int_{x_0}^{x}|s|t^{\Re(s)}2\varepsilon\,dt\\
        &\leq|s|\left[M(\varepsilon)\frac{x_0\Re(s)}{\Re(s)}+2\varepsilon\left(\frac{x^{\Re(s)}-x_0^{\Re(s)}}{\Re(s)}\right)\right]
    \end{align*}
    dividing by $|x^s|=x^{\Re(s)}$
    \begin{align*}
        \frac{|S(x)|}{x^{\Re(s)}}&\leq\frac{|s|}{\Re(s)}\left[M(\varepsilon)\left(\frac{x_0}{x}\right)^{\Re(s)}+2\varepsilon-2\varepsilon\left(\frac{x_0}{x}\right)^{\Re(s)}\right].
    \end{align*}    
    Due the dependance on $\varepsilon$ for the term on the right we can make this arbitrarily small. Hence, $\lim_{x\to\infty}\frac{S(x)}{x^s}=0$.
\end{proof}
Going back to studying the means of arithmetic functions:\\
Let 
\[
    M_f(x)=\frac{1}{x}\sum_{n\leq x}f(n)
\]
with $f:\N\to\C$, and 
\[
    M_f=\lim_{x\to\infty}M_f(x).
\]
\[
    L_f(x)=\frac{1}{\log x}\sum_{n\leq x}\frac{f(n)}{n},
\]
and 
\[
    L_f=\lim_{x\to\infty}L_f(x).
\]
\begin{theorem}{Existence of Logarithmic Means}{}
    Using the notation above, if $M_f$ exists then $L_f$ also exists and $M_f=L_f$.
\end{theorem}
\begin{remark}{}{}
    \[
        \frac{1}{x}\sum_{n\leq x}\Lambda(n)\to1\text{ as } x\to\infty.
    \]
    Also,
    \[
        \frac{1}{x}\sum_{n\leq x}\mu(n)\to0\text{ as } x\to\infty,
    \]
    \[
        M_\Lambda\Leftrightarrow \text{ PNT},
    \]
    computing
    \[
        M_\Lambda\& M_\mu\Rightarrow \text{ PNT}.
    \]
    However $L_\Lambda$ and $L_\mu$ are much more tractable objects and we will compute them in the next chapter.
\end{remark}
\begin{proof}[Proof of Theorem 2.13]
    Given $f:\N\to\C$ set 
    \[
        S_f(x)=\sum_{n\leq x}f(n)\quad\quad T_f(x)=\sum_{n\leq x}\frac{f(n)}{n}
    \]
    \[
        M_f(x)=\frac{S_f(x)}{x}\quad\quad L_f(x)=\frac{T_f(x)}{\log x}.
    \]
    Suppose that
    \[
        \lim_{x\to\infty}M_f(x)=M_f.
    \]
    We need to show that $\lim_{x\to\infty}L_f(x)=M_f$.
    \begin{align*}
        T(x)&=\sum_{n\leq x}\frac{f(n)}{n}\\
        &=\frac{S(x)}{x}+\int_{1}^{x}s(t)t^{-2}\,dt\\
        \frac{T(x)}{\log x}&=\frac{S(x)}{x\log x}+\frac{1}{\log x}\int_{1}^{x}\frac{S(t)}{t^2}\,dt\\
        &=\frac{S(x)}{x\log x}+\frac{1}{\log x}I(x)\\
        &=0+\frac{1}{\log x}I(x).
    \end{align*}
    Now we need only show that $\lim_{x\to\infty}\frac{I}{\log x}=M_f$. To do this we consider 
    \begin{align*}
        I(x)-M_f\log x&=\int_{1}^{x}\frac{S(t)}{t^2}\,dt-M_f\int_{1}^{x}\frac{1}{t}\,dt\\
        &=\int_{1}^{x}\frac{\frac{S(t)}{t}-M_f}{t}\,dt
    \end{align*}
    Thus
    \[
        |I(x)-M_f\log x|\leq\int_{1}^{x}\frac{\left|\frac{S(t)}{t}-M_f\right|}{t}\,dt.
    \]
    From here proceed as in Kronecker's lemma by splitting the integral into two intervals $[1,x_0]$ and $[x_0,x]$. Doing so will give the desired result.
\end{proof}
\addsec{Dirichlet Series}
Let $f:\N\to\C$, set 
\[
    F(s)=\sum_{n=1}^{\infty}\frac{f(n)}{n^s}.
\]
Then, $F(s)$ is the Dirichlet series associated with $f$.
\begin{theorem}{Convergence of Dirichlet Series}{}
    Let $f:\N\to\C$ be an arithmetic function, $F(s)$ the associated Dirichlet series and $s_f(x)=\sum_{n\leq x}f(n)$ the associated summatory function. Then,
    \begin{enumerate}[label=(\roman*)]
        \item For any $s\in\C$ with $\Re(s)>0$ such that $F(s)$ converges we have
        \begin{equation}\label{eqn:Dirichlet}
            F(s)=s\int_{1}^{\infty}\frac{s_f(x)}{x^{s+1}}\,dx.
        \end{equation}
        \item If $S_f(x)=\Oh(x^\alpha)$ for some $\alpha\geq0$, then $F(s)$ converges then for all $s$ with $\Re(s)>\alpha$ and equation \ref{eqn:Dirichlet} holds for all such $s$.
    \end{enumerate}
\end{theorem}

\begin{remark}
    \begin{enumerate}
        \item Let $\chi$ be a Dirichlet character ($|\chi(n)|=0$ or 1, $\chi$ is periodic).\\
        For non-trivial $\chi$, we can prove that
        \[
            \sum_{n\leq x}\chi(n)=\Oh(1).
        \]
        Let $L(s,\chi)$ be the Dirichlet series associated with $\chi$.
        \[
            L(s,\chi)=\sum_{n=1}^{\infty}\frac{\chi(n)}{n^s}.
        \]
        This is absolutly convergent for $\Re(s)>1$.\\
        By theorem 2.14(ii) and using the bound $\sum_{n\leq x}\chi(n)=\Oh(1)$ we see that $L(s,\chi)$ converges when $\Re(s)>0$.
        \item Theorem 2.14(i) can be expressed in terms of Mellin transforms.\\
        Given $\phi:\R^+\to\C$. The Mellin transform of $\phi$ is given by
        \[
            \overset{\sim}{\phi}(s)=\int_{0}^{\infty}\phi(x)x^{-s}\,dx
        \]
        ($s\in\C$) provided the integral converges.\\
        Theorem 2.14(i) can then be restated as $\dfrac{F(s)}{s}$ is the Mellin transform of $\dfrac{S_f(x)}{x}$.
    \end{enumerate}
\end{remark}
\begin{proof}[Proof of Theorem 2.14]
    Let $s\in\C$ be such that $\Re(s)>0$ and $F(s)$ converges.\\
    (i) Consdier 
    \[
        F_N(s)=\sum\frac{f(n)}{n^s}.
    \]
    By Abel's summation formula we have
    \begin{align*}
        F_N(s)&=\left(\sum_{n\leq N}f(n)\right)N^{-s}+s\int_{1}^{N}S_f(t)t^{-s-1}\,dt\\
        &=S_f(N)N^{-s}+s\int_{1}^{\infty}S_f(N)T^{-s-1}\,dt.
    \end{align*}
    As $N\to\infty$ $F_N(s)\to F(s)$, and by \textit{theorem 2.12} $\dfrac{S_f(N)}{N^s}\to0$.\\
    So as $N\to\infty$ we get
    \[
        F(s)=s\int_{1}^{\infty}S_f(t)t^{-s-1}\,dt.
    \]
    (ii) We have
    \[
        F_N(s)=\frac{S_f(N)}{N^s}+s\int_{1}^{N}S_f(t)t^{-s-1}\,dt.
    \]
    \[
        \left|\frac{S_f(N)}{N^s}\right|=\left|\frac{S_f(N)}{N^{\Re(s)}}\right|=\Oh\left(N^{\alpha-\Re(s)}\right),
    \]
    so for $\Re(s)>\alpha$, we have $\dfrac{S_f(N)}{N^s}\to 0$ as $N\to\infty$ we also have 
    \[
        \left|\frac{S_f(t)}{t^{s+1}}\right|=\Oh(t^{\alpha-\Re(s)-1}).
    \]
    It follows that
    \begin{align*}
        \left|\int_{1}^{N}\frac{S_f(t)}{t^{s+1}}\,dt\right|&\leq\int_{1}^{N}\left|\frac{S_f(t)}{t^{s+1}}\right|\,dt\\
        &=\Oh\left(\int_{1}^{N}t^{\alpha-\Re(s)-1}\,dt\right)\\
        &=\Oh\left(N^{\alpha-\Re(s)-1}-1\right).
    \end{align*}
    Hence, $\int_{1}^{\infty}\frac{S_f(t)}{t^{s-1}}\,dt$ is convergent for $\Re(s)>\alpha$. Taking $N\to\infty$ we get $F(s)$ is convergent for $\Re(s)>\alpha$ and 
    \[
        F(s)=s\int_{1}^{\infty}\frac{S_f(t)}{t^{s-1}}\,dt.\qedhere
    \]
\end{proof}
\addsec{Convolution Method}
Say we want to compute
\[
    \sum_{n\leq x}f(x)
\]
with $f:\N\to\C$, and suppose that $f=g\ast f_0$ where both $f_0,g$ are arithmetic functions.\\
We have
\begin{align*}
    \sum_{n\leq x}f(x)&=\sum_{n\leq x}g\ast f_0(n)\\
    &=\sum_{n\leq x}\sum_{d\mid n}g(d)f_0\left(\frac{d}{n}\right)\\
    &=\sum_{d\leq x}\sum_{\substack{n\leq x\\d\mid n}}g(d)f_0\left(\frac{d}{n}\right)\\
    &=\sum_{d\leq x}\sum_{m\leq \frac{x}{d}}g(d)f_0\left(m\right)\\
    &=\sum_{d\leq x}g(d)\sum_{m\leq \frac{x}{d}}f_0\left(m\right).
\end{align*}
If an asymptotic formula is known for the last sum we substitute and proceed.\\
Let us evaluate $\sum_{n\leq x}\phi(n)$.\\
We know that $\phi(n)=\mu\ast\mathrm{id}(n)$.
\begin{theorem}{Asymptotic Formula for Sum of $\phi(n)$}{}
    \[
        \sum_{n\leq x}\phi(n)=\frac{3}{\pi^2}x^2+\Oh\left(x\log x\right)
    \]
\end{theorem}
\begin{lemma}{Sum of $\frac{\mu(n)}{n^2}$}{}
    \[
        \sum_{n\leq x}\frac{\mu(n)}{n^2}=\frac{1}{\zeta(2)}.
    \]
    In general
    \[
        \sum_{n\leq x}\frac{\mu(n)}{n^s}=\frac{1}{\zeta(s)},\quad \Re(s)>1.
    \]
\end{lemma}
\begin{proof}
    \[
        1=\sum_{n=1}^{\infty}\frac{e(n)}{n^2}
    \]
    where 
    \[
        e(n)=
        \begin{cases}
            1&\text{if }n=1\\
            0&\text{otherwise}.    
        \end{cases}
    \]
    \begin{align*}
        1&=\sum_{n=1}^{\infty}\frac{e(n)}{n^2}\\
        &=\sum_{n=1}^{\infty}\frac{\displaystyle\sum_{d\mid n}\mu(d)}{n^2}\\
        &=\sum_{n=1}^{\infty}\frac{1}{n^2}\sum_{d\mid n}\mu(d)\\
        &=\sum_{d=1}^\infty\mu(d)\sum_{n=1}^{\infty}\frac{1}{n^2}\\
        &=\sum_{d=1}^\infty\mu(d)\sum_{m=1}^{\infty}\frac{1}{(md)^2}\\
        1&=\sum_{d=1}^\infty\frac{\mu(d)}{d^2}\zeta(2)\qedhere
    \end{align*}
\end{proof}
\begin{proof}[Proof of Theorem 2.15]
    We have
    \begin{align*}
        \sum_{n\leq x}\phi(n)&=\sum_{n\leq x}\mu\ast\mathrm{id}(n)\\
        &=\sum_{n\leq x}\sum_{d\mid n}\mu(d)\frac{d}{n}\\
        &=\sum_{d\leq x}\mu(d)\left(\sum_{m\leq\frac{n}{d}}m\right)\\
        &=\sum_{d\leq x}\mu(d)\left(\frac{\floor{\frac{x}{d}}\left(\floor{\frac{x}{d}+1}\right)}{2}\right).
    \end{align*}
    \begin{align*}
        \frac{1}{2}\floor{\frac{x}{d}}\left(\floor{\frac{x}{d}+1}\right)&=\frac{1}{2}\left(\frac{x}{d}-\fract{\frac{x}{d}}\right)\left(\frac{x}{d}-\fract{\frac{x}{d}}+1\right)\\
        &=\frac{1}{2}\left(\left(\frac{x}{d}\right)^2-2\fract{\frac{x}{d}}\frac{x}{d}+\frac{x}{d}-\fract{\frac{x}{d}}+\fract{\frac{x}{d}}^2\right)\\
        &=\frac{1}{2}\left(\frac{x}{d}\right)^2+\Oh\left(\frac{x}{d}\right).
    \end{align*}
    \begin{align*}
        \sum_{n\leq x}\phi(n)&=\sum_{d\leq x}\mu(d)\left(\frac{1}{2}\left(\frac{x}{d}\right)^2+\Oh\left(\frac{x}{d}\right)\right)\\
        &=\frac{x^2}{2}\sum_{d\leq x}\frac{\mu(d)}{d^2}+\Oh\left(\sum_{d\leq x}\frac{x}{d}\right)\\
        &=\frac{x^2}{2}\sum_{d\leq x}\frac{\mu(d)}{d^2}+\Oh(x\log x)\\
        &=\frac{x^2}{2}\left(\sum_{d=1}^{\infty}\frac{\mu(d)}{d^2}-\sum_{d> x}\frac{\mu(d)}{d^2}\right)+\Oh(x\log x)\\
        &=\frac{x^2}{2\zeta(2)}-\frac{x^2}{2}\sum_{d> x}\frac{\mu(d)}{d^2}+\Oh(x\log x).
    \end{align*}
    \begin{align*}
        \left|\sum_{d>x}\frac{\mu(d)}{d^2}\right|&\leq\sum_{d\geq\floor{x}+1}\frac{\mu(d)}{d^2}\\
        &=\int_{\floor{x}}^{\infty}\frac{dt}{t^2}\\
        &=\Oh\left(\frac{1}{x}\right).
    \end{align*}
    All in all we have
    \begin{align*}
        \sum_{n\leq x}\phi(n)&=\frac{x^2}{2\zeta(2)}+\Oh\left(x\log x\right)\\
        &=\frac{3}{\pi^2}x^2+\Oh\left(x\log x\right).\qedhere
    \end{align*}
\end{proof}
\addsec{Dirichlet's Hyperbola Method}
Suppose we have two arithmetic functions
\[
    f,g:\N\to\C.
\]
Let
\[
    F(x)=\sum_{n\leq x}f(n),\quad G(x)=\sum_{n\leq x}g(x).
\]
Consider $h=f\ast g$.\\
i.e.,
\[
    h(n)=\sum_{d\mid n}f(d)g\left(\frac{n}{d}\right).
\]
We have
\begin{align*}
    H(x)=\sum_{n\leq x}h(n)&=\sum_{n\leq x}\sum_{d\mid n}f(d)g\left(\frac{n}{d}\right)\\
    &=\sum_{n\leq x}\sum_{n=ab}f(a)g(b)\\
    &=\sum_{\substack{a,b\\ ab\leq x}}f(a)g(b)
    \intertext{let $y>0$}
    &=\sum_{\substack{ab\leq x\\ a\leq y}}f(a)g(b)+\sum_{\substack{ab\leq x\\ a> y}}f(a)g(b)\\
    &=\sum_{a\leq y}f(a)\sum_{b\leq\frac{x}{a}}g(b)+\sum_{b\leq \frac{x}{y}}g(b)\sum_{y<a\leq\frac{x}{b}}f(a)\\
    &=\sum_{a\leq y}f(a)G(\frac{x}{a})+\sum_{b\leq \frac{x}{y}}g(b)\left(F\left(\frac{x}{b}\right)-F(y)\right)\\
    &=\sum_{a\leq y}f(a)G(\frac{x}{a})+\sum_{b\leq \frac{x}{y}}g(b)F\left(\frac{x}{b}\right)-\sum_{b\leq \frac{x}{y}}g(b)F(y)\\
    &=\sum_{a\leq y}f(a)G\left(\frac{x}{a}\right)+\sum_{b\leq \frac{x}{y}}g(b)F\left(\frac{x}{b}\right)-G\left(\frac{x}{y}\right)F(y)
\end{align*}
If we choose $y=\sqrt{x}$ (this choice depends on the specific functions you are dealing with), the above can be written as 
\[
    \sum_{n\leq x}f\ast g(n)=\sum_{a\leq \sqrt{x}}f(a)G\left(\frac{x}{a}\right)+\sum_{b\leq \sqrt{x}}g(b)F\left(\frac{x}{b}\right)-G(\sqrt{x})F(\sqrt{x}).
\]
Why is this called the ``hyperbola method''? If we use this method for $\sum_{n\leq x}d(n)$ we have
\[
    \sum_{n\leq x}d(n)=\sum_{n\leq x}\sum_{ab=n}1=\sum_{\substack{a,b\\ab\leq x}}1.
\]
We can see that the first sum, sums over lattice points under hyperbola $ab=x$ of points with $a\leq\sqrt{x}$ and the second sums over points with $b\leq\sqrt{x}$ and the thid term subtracts one set of points that have been double counted.
\begin{theorem}{Sum of Divisor Function}{}
    \[
        \sum_{n\leq x}\mathrm{d}(n)=x\log x+(2\gamma-1)x+\Oh(\sqrt{x}).
    \]
\end{theorem}
\begin{proof}[Proof of Theorem 2.17]
    \begin{align*}
        \sum_{n\leq x}d(n)=\sum_{\substack{a,b\\ ab\leq x}}1&=\sum_{a\leq \sqrt{x}}\sum_{b\leq \frac{x}{a}}1+\sum_{b\leq \sqrt{x}}\sum_{a\leq \frac{x}{b}}1-\sum_{b\leq \sqrt{x}}\sum_{b\leq \sqrt{x}}1\\
        &=\sum_{a\leq \sqrt{x}}\floor{\frac{x}{a}}+\sum_{b\leq \sqrt{x}}\floor{\frac{x}{b}}-\sum_{b\leq \sqrt{x}}\floor{\sqrt{x}}\\
        &=2\sum_{a\leq \sqrt{x}}\floor{\frac{x}{a}}-\floor{\sqrt{x}}\floor{\sqrt{x}}\\
        &=2\sum_{a\leq \sqrt{x}}\left(\frac{x}{a}-\fract{\frac{x}{a}}\right)-(\sqrt{x}-\fract{\sqrt{x}})^2\\
        &=2x\sum_{a\leq \sqrt{x}}\frac{1}{a}-2\sum_{a\leq \sqrt{x}}\fract{\frac{x}{a}}-x-2\sqrt{x}\fract{\sqrt{x}}-\fract{\sqrt{x}}^2\\
        &=2x\left(\log\sqrt{x}+\gamma+\Oh\left(\frac{1}{\sqrt{x}}\right)\right)-x+\Oh(\sqrt{x})\\
        &=2\log x+(2\gamma-1)x+\Oh(\sqrt{x}).\qedhere
    \end{align*}
\end{proof}
\begin{remark}
    Dirichlet's Divisor Problem\\
    \[
        \Delta(x)=D(x)-(x\log x+(2\gamma-1)x)\tag{$\ast$}
    \]
    \[
        \text{Theorem 2.17}\Rightarrow \Delta(x)=\Oh(\sqrt{x})
    \]
    \[
        \text{Conjecture:} \Delta(x)=\Oh(x^{\frac{1}{4}}).
    \]
    Hardy et al. showed that $\Delta(x)=\Oh(x^{\frac{1}{3}})$. The state of the art is $\Delta(x)=\Oh(x^\theta)$ with $\theta=0.31\dots$.\\
    $d(n)\leq2$, $d(n)\geq\exp\left(\frac{(1+\varepsilon)\log2\log n}{\log\log n}\right)$ for any $\varepsilon>0$. $d(n)\leq C(\varepsilon)n^\varepsilon$ for any $\varepsilon>0$. $d(n)\ll_\varepsilon n^\varepsilon$.
\end{remark}
\chapter{Elementary Results for the \\Prime Coutning Function}
The PNT 
\[
    \pi(x)\sim\frac{x}{\log x}
\]
was proven independantly by Hadamard and de la Vall\'{e} Poussin in the 1890's. Chebyshev proved (1848-1850) the following in equality: $\exists c_1,c_2>0$ such that
\[
    c_1\frac{x}{\log x}\leq \pi(x)\leq c_2\frac{x}{\log x}    
\]
for all $x\geq 2$, in other words
\[
    \pi(x)\asymp\frac{x}{\log x}.
\]
This inequality is enough to prove Bertrand's postulate (there always exists a prime between $n$ and $2n$).\\
Chebyshev's functions are given by
\[
    \Psi(x)=\sum_{n\leq x}\Lambda(n),\quad \theta(x)=\sum_{\substack{p\leq x\\p \text{ prime}}}\log p.
\]
\begin{theorem}{Asymptotic Equivalence for Chebyshev's Functions}{}
    For all $x\geq2$, we have:
    \begin{enumerate}
        \item $\Psi(x)\asymp x$, i.e., $\exists c_1,c_2>0$ such that $c_1x\leq \Psi(x)\leq c_2x$ for all $x\geq2$.
        \item $\theta(x)\asymp x$
        \item $\pi(x)\asymp\frac{x}{\log x}$
    \end{enumerate}
\end{theorem}
\begin{proof}[Proof of Theorem 3.1]
    We will prove (1) and then deduce (2) and (3).\\
    Consider
    \begin{align*}
        S(x)&=\sum_{n\leq x}\log n\\
        &=\sum_{n\leq x}\sum_{d\mid n}\Lambda(d)\\
        &=\sum_{d\leq x}\sum_{m\leq\frac{x}{d}}\Lambda(d)\\
        &=\sum_{d\leq x}\Lambda(d)\sum_{m\leq\frac{x}{d}}1\\
        S(x)&=\sum_{d\leq x}\Lambda(d)\floor{\frac{x}{d}}.
    \end{align*}
    Now we introduce
    \begin{align*}
        D(x)&=S(x)-2S\left(\frac{x}{2}\right)\\
        &=\sum_{d\leq x}\Lambda(d)\floor{\frac{x}{d}}-2\sum_{d\leq\frac{x}{2}}\Lambda(d)\floor{\frac{x}{2d}}\\
        &=\sum_{d\leq x}\Lambda(d)\floor{\frac{x}{d}}-2\sum_{d\leq x}\Lambda(d)\floor{\frac{x}{2d}}+\sum_{\frac{x}{2}<d\leq x}\Lambda(d)\floor{\frac{x}{2d}}\\
        &=\sum_{d\leq x}\Lambda(d)\left(\underbrace{\floor{\frac{x}{d}}-2\floor{\frac{x}{2d}}}_{f(x)}\right)
        \intertext{where 
        \(f(t)=\floor{t}-2\floor{\frac{t}{2}}=
        \begin{cases}
            1, & 1\leq t\leq 2\\
            0, & t\geq 2
        \end{cases}\)}
        &=\sum_{d\leq x}\Lambda(d)f\left(\frac{x}{d}\right)\\
        &=\underbrace{\sum_{d\leq\frac{x}{2}}\Lambda(d)f\left(\frac{x}{d}\right)}_{\leq0}+\sum_{\frac{x}{2}<d\leq x}\Lambda(d)f\left(\frac{x}{d}\right)\\
        &\geq\sum_{\frac{x}{2}<d\leq x}\Lambda(d)f\left(\frac{x}{2}\right)\\
        &=\sum_{\frac{x}{2}<d\leq x}\Lambda(d)\\
        &=\Psi(x)-\Psi\left(\frac{x}{2}\right).
    \end{align*}
    Here, $D(x)\geq \Psi(x)-\Psi\left(\frac{x}{2}\right)$.\\
    On the other hand
    \begin{align*}
        D(x)&=S(x)-2S\left(\frac{x}{2}\right)\\
        &=\sum_{n\leq x}\log n-2\sum_{n\leq\frac{x}{2}}\log n\\
        &=x(\log x-1)+\Oh(\log x)-2\frac{x}{2}\left(\log\frac{x}{2}-1\right)+\Oh(\log x)\\
        D(x)&=(\log2)x+\Oh(\log x).
    \end{align*}
    This means that there exists $A>0$ such that $D(x)=(\log2)x+g(x),$ where $|g(x)|\leq A\log x$ for all $x$ sufficiently large. And so $D(x)\leq(\log2)x+A\log x$ for all $x$ sufficiently large. Let $\epsilon\geq0$ be given for $x>N$ (for some $N>0$) we have 
    \[
        \log x\leq \frac{\epsilon}{A}x
    \]
    and so for $x$ large enough say $x>x_0$ we have
    \[
        D(x)\leq(\log2+\epsilon)x
    \]
    since $g(x)>-A\log x$ for $x$ sufficiently large. We also have
    \begin{align*}
        D(x)&\geq(\log2)x-A\log x\\
        &\geq(\log2-\epsilon)x.
    \end{align*}
Therefore, 
\[
    (\log2-\epsilon)x\leq D(x)\leq(\log2+\epsilon)x
\]
for all $x\geq x_0$. We also have,
\[
    D(x)\geq\Psi(x)-\Psi\left(\frac{x}{2}\right)
\]
so we have
\[
    \Psi(x)-\Psi\left(\frac{x}{2}\right)\leq D(x)\leq\Psi(x).
\]
In particular, we get
\[
    \Psi(x)\geq(\log2-\epsilon)x
\]
for $x\geq x_0$. This gives us 
\[
    \Psi(x)\geq D(x)\geq(\log2-\varepsilon)x
\]
for all $x\geq x_0$. We also have that 
\begin{align*}
    \Psi(x)&\leq D(x)+\Psi\left(\frac{x}{2}\right)\\
    &\leq(\log 2-\varepsilon)x+\Psi\left(\frac{x}{2}\right)& \forall x\geq x_0
    \intertext{By choosing a small enough $\varepsilon$}
    &\leq x+\Psi\left(\frac{x}{2}\right)& \forall x\geq x_0\\
    &\leq x+\frac{x}{2}+\Psi\left(\frac{x}{4}\right)& \forall \frac{x}{2}\geq x_0\\
    &\quad\quad\quad\quad\:\,\vdots\\
    &\leq x+\frac{x}{2}+\frac{x}{4}+\cdots+\Psi\left(\frac{x}{2^k}\right)& \forall \frac{x}{2^{k-1}}\geq x_0.
\end{align*}
Choose $k$ to be the largest integer such that 
\[
    \frac{x}{2^{k-1}}\geq x_0
\]
i.e.,
\[
    \frac{x}{2^k}\leq x_0\leq\frac{x}{2^{k-1}}.
\]
Since $\Psi$ is an increasing function, we have 
\[
    \Psi(x_0)\geq\Psi\left(\frac{x}{2}\right).
\]
It follows that 
\begin{align*}
    \Psi(x)&\leq x+\frac{x}{2}+\cdots+\frac{x}{2^{k-1}}+\Psi(x_0)\\
    &=x\sum_{i=0}^{k-1}\left(\frac{1}{2}\right)^i+\Psi(x_0)\\
    &=2x+\Psi(x_0)
    \intertext{For large enough $x$ we have $\Psi(x_0)\leq \varepsilon x$}
    \Psi(x)&\leq(2+\varepsilon)x
    \intertext{For all $x\geq X_0$, and some $X_0>1$.}
\end{align*}
Combining everything, we get
\[
    (2-\varepsilon)\leq\Psi(x)\leq(2+\varepsilon)x
\]
for all $x\geq X_0$.\\
What if $2\leq x\leq X_0$? Trivially we have that 
\[
    \frac{\Psi(x)}{x}\leq\frac{\Psi(X_0)}{2}.
\]
We also have that 
\[
    \frac{\Psi(x)}{x}\geq\frac{\Psi(2)}{x}=\frac{\log2}{x}\leq\frac{\log2}{X_0}.
\]
This gives
\[
    \frac{\log2}{X_0}x\leq\Psi(x)\leq\frac{\Psi(X_0)}{2}x
\]
for all $2\leq x\leq X_0$. All in all, $\exists c_1,c_2>0$ such that 
\[
    c_1x\leq \Psi(x)\leq c_2x
\]
for all $x\geq2$.\\
Now for part (2)
\begin{align*}
    \Psi(x)-\theta(x)&=\sum_{n\leq x}\Lambda(n)-\sum_{p\leq x}\log p\\
    &=\sum_{\substack{n\leq x\\ n=p^m}}\log p-\sum_{p\leq x}\log p\\    
    &=\sum_{\substack{p^m\leq x\\ p\text{ prime}\\m\in\N}}\log p-\sum_{p\leq x}\log p\\
    &=\sum_{\substack{p^m\leq x\\ p\text{ prime}\\m\geq2}}\log p\\
    &=\sum_{p\leq\sqrt{x}}\sum_{\substack{m\leq\frac{\log x}{\log p}\\m\geq2}}\log p\\    
    &=\sum_{p\leq\sqrt{x}}\log p\floor{\frac{\log x}{\log p}}\\    
    &\leq\sum_{p\leq\sqrt{x}}\log p\frac{\log x}{\log p}\\
    &=\log x\sum_{p\leq\sqrt{x}}1\\
    &\leq(\log x)\sqrt{x}.
\end{align*}
We just showed that
\[
    \Psi(x)-\theta(x)\leq (\log x)\sqrt{x}.
\]
Now we have
\[
    \theta(x)\leq \Psi(x)\leq c_2x
\]
we also have
\begin{align*}
    \theta(x)&\geq\Psi(x)-(\log x)\sqrt{x}\\
    &\geq c_1x-(\log x)\sqrt{x}
\end{align*}
for $x$ large enough we have
\[
    (\log x)\sqrt{x}\leq\frac{c_1}{2}x
\]
and so 
\[
    \theta(x)\geq c_1x-\frac{c_1}{2}x=\frac{c_1}{2}x.
\]
\\
Now for part (3)
\begin{align*}
    \pi(x)&=\sum_{p\leq x}1\\
    &\geq\sum_{p\leq x}\frac{\log p}{\log x}\\
    &=\frac{1}{\log x}\sum_{p\leq x}\log p\\
    &=\frac{1}{\log x}\theta(x)\geq\frac{\frac{c_1}{x}x}{\log x}\\
    &=\sum_{p\leq\sqrt{x}}1+\sum_{\sqrt{x}<p\leq x}1\\
    &\leq\sqrt{x}+\frac{1}{\log\sqrt{x}}\sum_{\sqrt{x}<p\leq x}\log\sqrt{x}\\
    &\sqrt{x}+\frac{2}{\log x}\sum_{\sqrt{}\leq p\leq x}\log p\\
    &\leq\sqrt{x}+\frac{2}{\log x}(\theta(x)-\theta(\sqrt{x}))\\
    &\leq\sqrt{x}+\frac{2}{\log x}(c_2x-\frac{c_1}{2}\sqrt{x})\\
    &\leq C\frac{x}{\log x}
\end{align*}
for some $C>0$. Indeed we have that 
\[
    \pi(x)\asymp\frac{x}{\log x}
\]
for $x$ large enough. Repeating the argument we did to bound $\Psi(x)$ for small values of $x$, we get $\pi(x)\asymp\frac{x}{\log x}$ for all $x\geq 2$.\qedhere
\end{proof}    
\begin{theorem}{}{}
    For all $x\leq2$, we have 
    \begin{enumerate}
        \item[(1)] $\theta(x)=\Psi(x)+\Oh(\sqrt{x})$;
        \item[(2)] $\pi(x)=\frac{\Psi(x)}{\log x}+\Oh\left(\frac{x}{\log^2x}\right)$.
    \end{enumerate}
\end{theorem}
\begin{proof}[Proof of Theorem 3.2]
    \begin{enumerate}
        \item[(1)]
        \begin{align*}
            \Psi(x)-\theta(x)&=\vdots\\
            &\leq\log x\sum_{p\leq\sqrt{x}}1\\
            &=\log x\pi(\sqrt{x})
            \intertext{By theorem 3.1 part 3}
            &\leq\log x\frac{c_2\sqrt{x}}{\log\sqrt{x}}\\
            &\leq2c_2\sqrt{x}\\
            &\ll\sqrt{x}.
        \end{align*}
        \item[(2)]
        \begin{align*}
            \pi(x)&=\sum_{p\leq x}1\\
            &=\sum_{p\leq x}\frac{\log p}{\log p}\\
            &=\left(\sum_{p\leq x}\log p\right)\frac{1}{\log x}+\int_{2}^{x}\left(\sum_{p\leq t}\log p\right)\frac{1}{t(\log t)^2}\,dt\\
            &=\frac{\Oh(x)}{\log x}+\int_{2}^{x}\frac{\Oh(t)}{t(\log t)^2}\,dt.
        \end{align*}
        Observe that
        \[
            \int_{2}^{x}\frac{\Oh(t)}{t(\log t)^2}\,dt\ll\int_{2}^{x}\frac{1}{(\log t)^2}\,dt=\Li_2(x)\ll\frac{x}{(\log x)^2}.
        \]
        Hence,
        \begin{align*}
            \pi(x)&=\frac{\Oh(x)}{\log x}+\Oh\left(\frac{x}{(\log x)^2}\right)\\
            &=\frac{\Psi(x)}{\log x}+\Oh\left(\frac{\sqrt{x}}{\log x}\right)+\Oh\left(\frac{x}{(\log x)^2}\right)\\
            &=\frac{\Psi(x)}{\log x}+\Oh\left(\frac{x}{(\log x)^2}\right).\qedhere
        \end{align*}   
    \end{enumerate}
\end{proof}
\begin{corollary}{}{}
    The following are all equivalent as $x\to\infty$
    \begin{enumerate}
        \item[(1)] $\pi(x)\sim\frac{x}{\log x}$;
        \item[(2)] $\theta(x)\sim x$;
        \item[(3)] $\Psi(x)\sim x$.   
    \end{enumerate}
\end{corollary}
\addsec{Mertens' Estimate}
\begin{theorem}{}{}
    We have
    \begin{enumerate}
        \item[(1)] 
        \[
            \sum_{n\leq x}\frac{\Lambda(n)}{n}=\log x+\Oh(1)
        \]
        \item[(2)]
        \[
            \sum_{p\leq x}\frac{\log p}{p}=\log x+\Oh(1)
        \]
        \item[(3)]
        \[
            \sum_{p\leq x}\frac{1}{p}=\log\log x+A+\Oh\left(\frac{1}{\log x}\right)
        \]
        \item[(4)] Mertens' formula
        \[
            \prod_{p\leq x}\left(1-\frac{1}{p}\right)=\frac{e^{-\gamma}}{\log x}\left(1+\Oh\left(\frac{1}{\log x}\right)\right).
        \]
    \end{enumerate}
\end{theorem}
\begin{proof}[Proof of Theorem 3.4]
    \begin{enumerate}
        \item[(1)] Let 
        \begin{align*}
            S(x)&=\sum_{n\leq x}\log n\\
            &=x\log x-x+\Oh(\log x),
        \end{align*}
        on the other hand we have
        \begin{align*}
            S(x)&=\sum_{d\leq x}\Lambda(d)\floor{\frac{x}{d}}\\
            &=\sum_{d\leq x}\Lambda(d)\left(\frac{x}{d}-\fract{\frac{x}{d}}\right)\\
            &=x\sum_{d\leq x}\frac{\Lambda(d)}{d}+\sum_{d\leq x}\Lambda(d)\fract{\frac{x}{d}}\\
            &=x\sum_{d\leq x}\frac{\Lambda(d)}{d}+\Oh\left(\sum_{d\leq x}\Lambda(d)\right)\\
            &=x\sum_{d\leq x}\frac{\Lambda(d)}{d}+\Oh(x).
        \end{align*}
        Since $S(x)=x\log x+\Oh(x)$ we get
        \[
            \log x+\Oh(x)=\sum_{d\leq x}\frac{\Lambda(d)}{d}+\Oh(x).
        \]
        \item[(2)]
        \begin{align*}
            \sum_{p\leq x}\frac{\log p}{p}&=\sum_{p^m\leq x}\frac{\log p}{p^m}-\sum_{\substack{p^m\leq x\\ m\leq2}}\frac{\log p}{p^m}\\
            &=\sum_{n\leq x}\frac{\Lambda(n)}{n}-\sum_{\substack{p^m\leq x\\ m\leq2}}\frac{\log p}{p^m}.
        \end{align*}
        Note that
        \begin{align*}
            -\sum_{\substack{p^m\leq x\\ m\leq2}}\frac{\log p}{p^m}&=\sum_{p\leq\sqrt{x}}\log p\sum_{2\leq m\leq\frac{\log x}{\log p}}\frac{1}{p^m}\\
            &\leq\sum_{p}\log p\sum_{m=1}^{\infty}\left(\frac{1}{p}\right)^m\\
            &=\sum_{p}\log p\left(\frac{p^{-2}}{1-p^{-1}}\right)\\
            &=\sum_{p}\frac{\log p}{p^2-p}<\infty\\
            &=\Oh(1).
        \end{align*}
        \item[(3)] 
        \begin{align*}
            \sum_{p\leq x}\frac{1}{p}&=\sum_{p\leq x}\frac{1}{p}\frac{\log p}{\log p}\\
            &=\left(\sum_{p\leq x}\frac{\log p}{p}\right)\frac{1}{\log x}+\int_{2}^{x}\left(\sum_{p\leq t}\frac{\log p}{p}\right)\frac{1}{t(\log t)^2}\,dt 
            \intertext{Let $R(x)=\Oh(1)$}
            &=\frac{\log p+R(x)}{\log x}+\int_{2}^{x}\frac{\log t+R(t)}{t(\log t)^2}\,dt\\
            &=1+\frac{R(x)}{\log x}+\int_{2}^{x}\frac{1}{t\log t}\,dt+\int_{2}^{x}\frac{R(x)}{t(\log t)^2}\,dt\\
            &=1+\log\log x-\log\log2+\int_{2}^{x}\frac{R(t)}{t(\log t)^2}\,dt+\Oh\left(\frac{1}{\log x}\right).
        \end{align*}
        Let's consider $\int_{2}^{x}\frac{R(t)}{t(\log t)^2}\,dt$. We know that $\int_{2}^{\infty}\frac{R(t)}{t(\log t)^2}\,dt$ is convergent, since $\int_{2}^{\infty}\frac{R(t)}{t(\log t)^2}\,dt\ll\int_{\infty}^{x}\frac{1}{t(\log t)^2}\,dt<\infty$. Let's call
        \[
            \int_{2}^{\infty}\frac{R(t)}{t(\log t)^2}\,dt=C.
        \]
        Now
        \[
            \int_{2}^{\infty}\frac{R(t)}{t(\log t)^2}\,dt\ll\int_{2}^{\infty}\frac{1}{t(\log t)^2}\,dt=\Oh\left(\frac{1}{\log x}\right).
        \]
        Hence, 
        \[
            \sum_{p\leq x}\frac{1}{p}=\log\log x+1-\log\log2+C+\Oh\left(\frac{1}{\log x}\right).
        \]
        \item[(4)] We need to show that 
        \[
            \prod_{p\leq x}\left(1-\frac{1}{p}\right)=\frac{e^{-\gamma}}{\log x}\left(1+\Oh\left(\frac{1}{\log x}\right)\right).
        \]
        For now we will show that 
        \[
            \prod_{p\leq x}\left(1-\frac{1}{p}\right)=\frac{C}{\log x}\left(1+\Oh\left(\frac{1}{\log x}\right)\right)
        \]
        for some $C>0$. In the next chapter we will establish that $C=e^{-\gamma}$.\\
        Our problem is reduced to proving
        \[
            \prod_{p\leq x}\left(1-\frac{1}{p}\right)=\log C-\log\log x+\log\left(1+\Oh\left(\frac{1}{\log x}\right)\right).
        \]
        Notice that $|\log(1+y)|\asymp |y|$ for $|y|\leq\frac{1}{2}$. So for $x$ large enough the term $\Oh\left(\frac{1}{x}\right)$ can be made less than $\frac{1}{2}$, and so 
        \[
            \log\left(1+\log\left(\frac{1}{\log x}\right)\right)=\Oh\left(\frac{1}{\log x}\right).
        \]
        Let's prove
        \[
            -\sum_{p\leq x}\log\left(1-\frac{1}{p}\right)=\log\log x+\log C+\Oh\left(\frac{1}{\log x}\right).
        \]
        Note that 
        \[
            \log(1-x)=-\sum_{n=1}^{\infty}\frac{x^n}{n}.
        \]
        So, we have
        \begin{align*}
            -\sum_{p\leq x}\log\left(1-\frac{1}{p}\right)&=\sum_{p\leq x}\sum_{n=1}^{\infty}\frac{\left(\frac{1}{p}\right)^n}{n}\\
            &=\sum_{p\leq x}\sum_{n=1}^{\infty}\frac{1}{np^n}\\
            &=\sum_{p\leq x}\left(\frac{1}{p}+\sum_{n=2}^{\infty}\frac{1}{np^n}\right)\\
            &=\sum_{p\leq x}\frac{1}{p}+\sum_{p\leq x}\sum_{n=2}^{\infty}\frac{1}{np^n}.
        \end{align*}
        Consider the second sum
        \[
            \sum_{p\leq x}\underbrace{\frac{n=2}{\infty}\frac{1}{np^n}}_{r_p}.
        \]
        We have
        \begin{align*}
            r_p=\sum_{n=2}^{\infty}&\leq\frac{1}{2}\sum_{n=2}^{\infty}\frac{1}{p^n}\\
            &=\frac{1}{2}\left(\frac{\frac{1}{p^2}}{1-\frac{1}{p}}\right)\\
            &=\frac{1}{2}\left(\frac{1}{p^2-p}\right).
        \end{align*}
        Now 
        \[
            \sum_{p\leq x}r_p=\sum_{p=2}^{\infty}r_p-\sum_{p>x}r_p
        \]
        since $r_p\leq\frac{1}{2}\left(\frac{1}{p^2-p}\right)$, the comparison test implies that $\sum_{p=2}^{\infty}r_p$ is convergent to say $B$. We also have, 
        \begin{align*}
            \sum_{p>x}r_p&\leq\frac{1}{2}\sum_{p>x}\frac{1}{p^2-p}\\
            &\leq\frac{1}{2}\sum_{n>x}\frac{1}{n^2-n}\\
            &\leq\sum_{n\geq x}\frac{1}{n^2}\\
            &\ll\int_{x}^{\infty}\frac{1}{t^2}\,dt\\
            &=\Oh\left(\frac{1}{x}\right).
        \end{align*}
        Therefore,
        \[
            \sum_{p\leq x}\sum_{n=2}^{\infty}\frac{1}{np^n}=B+\Oh\left(\frac{1}{x}\right).
        \]
        It follows that
        \[
            -\sum_{p\leq x}\log\left(1-\frac{1}{p}\right)=\log\log x+A+B+\Oh\left(\frac{1}{\log x}\right).\qedhere
        \]     
    \end{enumerate}
\end{proof}
\begin{corollary}{}{}
    We have
    \[
        \int_{1}^{x}\frac{\Psi(t)}{t^2}\,dt=\log x+\Oh(1).
    \]    
\end{corollary}
\begin{proof}
    \begin{align*}
        \log x+\Oh(1)&=\sum_{n\leq x}\frac{\Lambda(n)}{n}\\
        &=\left(\sum_{n\leq x}\Lambda(n)\right)\frac{1}{x}+\int_{1}^{x}\left(\sum_{n\leq t}\Lambda(n)\right)\frac{1}{t^2}\,dt\\
        &=\frac{\Psi(x)}{x}+\int_{1}^{x}\frac{\Psi(t)}{t^2}\,dt.
    \end{align*}
    Hence, 
    \[
        \int_{1}^{x}\frac{\Psi(t)}{t^2}\,dt=\log x+\Oh(1)
    \]
    which is the desired result.
\end{proof}

\addsec{Consequences of the PNT}
\begin{theorem}{}{}
    The PNT implies the following
    \begin{enumerate}
        \item[(1)] The $n^{\text{th}}$ prime satisfies $p_n\sim n\log n$ as $n\to\infty$;
        \item[(2)] $\omega(n)=\#$ of distinct prime factors of $n$,
        \[
            \limsup_{n\to\infty}\frac{\omega(n)}{\frac{\log n}{\log\log n}}=1;
        \]
        \item[(3)] Given $\varepsilon>0$, $\exists x_0=x_0(\varepsilon)$ such that for all $x\geq x_0$ there exists a prime $p$ with
        \[
            x<p\leq(1+\varepsilon)x;
        \]
        \item[(4)] The set 
        \[
            \left\{\frac{p}{q}:p,q\text{ prime}\right\}
        \]
        is dense in $\R_+$;
        \item[(5)] $a_i\in\{1,2,3,\dots,9\}$. Consider a string $a_1,a_2,\dots,a_n$ $(a_i\neq0)$, then there exists a prime $\#$ whose decimal expansion starts with $a_1a_2\dots a_n$.
    \end{enumerate}
\end{theorem}
\begin{proof}[Proof of Theorem 3.6]
    \begin{enumerate}
        \item[(1)]
        \[
            \text{PNT}\Rightarrow n=\pi(p_n)\sim\frac{p_n}{\log p_n}.
        \]
        Let's show that $\log n\sim\log p_n$ that is $\lim_{n\to\infty}\frac{\log n}{\log p_n}=1$. We have
        \[
            \lim_{n\to\infty}\frac{n}{\frac{p_n}{\log p_n}}=1
        \]
        given $\varepsilon>0$, $\exists N\geq1$ such that
        \[
            \left|\frac{n}{\frac{p_n}{\log p_n}}-1\right|\leq\varepsilon
        \]
        for all $n\geq N$.
        \begin{align*}
            1-\varepsilon&\leq\frac{n}{\frac{p_n}{\log p_n}}&\leq1+\varepsilon\\
            \frac{p_n}{\log p_n}(1-\varepsilon)&\leq n&\leq\frac{p_n}{\log p_n}(1+\varepsilon)\\
            \log p_n-\log\log p_n+\log(1-\varepsilon)&\leq\log n&\leq\log p_n-\log\log p_n+\log(1+\varepsilon)\\
            1-\frac{\log\log p_n}{\log p_n}+\frac{\log(1-\varepsilon)}{\log p_n}&\leq\frac{\log n}{\log p_n}&\leq1-\frac{\log\log p_n}{\log p_n}+\frac{\log(1+\varepsilon)}{\log p_n}.
        \end{align*}
        Hence we have
        \[
            \lim_{n\to\infty}\frac{\log n}{\log p_n}=1.\qedhere
        \]
    \end{enumerate}
\end{proof}
For the rest of the chapter, we will prove the following theorem.
\begin{theorem}{}{}
    \[
        PNT\Leftrightarrow\frac{1}{x}\sum_{n\leq x}\mu(n)\to0
    \]
    as $x\to\infty$. That is $\sum_{n\leq x}\mu(n)=\oh(x)$.
\end{theorem}
\begin{remark}
    Let 
    \[
        Q(x)=\sum_{n\leq x}\mu^2(n)=\# \text{ of square free integers up to } x,
    \]
    \[
        Q_-(x)=\sum_{n\leq x}\mu^2(n)=\# \text{ of square free integers up to } x \text{ with an odd number of prime divisors},
    \]
    \[
        Q_+(x)=\sum_{n\leq x}\mu^2(n)=\# \text{ of square free integers up to } x \text{ with an even number of prime divisors}.
    \]
    See that
    \begin{align*}
        Q(x)&=Q_-(x)+Q_+(x)
        &=\frac{6}{\pi^2}x+\Oh(\sqrt{x}).
    \end{align*}
    We have 
    \[
        \frac{1}{x}\sum_{n\leq x}\mu(n)=\frac{1}{x}\left(Q_+(x)-Q_-(x)\right).
    \]
    If 
    \[
        \frac{1}{x}\sum_{n\leq x}\mu(n)\to0
    \]
    as $x\to\infty$, then
    \[
        Q_-(x)=Q_+(x)+\oh(x)=\frac{3}{\pi^2}x+\oh(x).
    \]
    So the statement of theorem 3.7 is telling us that a statement about the distribution of prime numbers si equivalent to a statement about the distribution of squarefree numbers with an odd number of prime factors.
\end{remark}
In preparing for the proof of theorem 3.7, we require the following lemmas.
\begin{lemma}{}{}
    \[
        \left|\sum_{n\leq x}\frac{\mu(n)}{n}\right|\leq1
    \]
\end{lemma}
\begin{corollary}{}{}
    Let 
    \[
        L(\mu)=\lim_{x\to\infty}\frac{1}{\log x}\sum_{n\leq x}\frac{\mu(n)}{n}.
    \]
    Then $L(\mu)=0$, and if $M(\mu)$ exists, then $M(\mu)=0$. 
\end{corollary}
\begin{lemma}{}{}
    $f:\N\to\C$
    \[
        H(f)=\lim_{x\to\infty}\frac{1}{x\log x}\sum_{n\leq x}f(n)\log n.    
    \]
    $H(f)$ exists \textit{if and only if} $M(f)$ exists.
\end{lemma}
Recall that the goal is to prove theorem 3.7
\[
    \text{PNT}\Leftrightarrow\sum_{n\leq x}\mu(n)=\oh(x)
\]
i.e.,
\[
    \lim_{n\to\infty}\frac{1}{x}\sum_{n\leq x}\mu(n)=0.
\]
To this end we require lemma 3.7, 3.8 and 3.11.
\begin{lemma}{}{}
    \[
        \mu(n)\log n=-(\mu\ast\Lambda)(n),\quad n\in\N
    \]
\end{lemma}
\begin{proof}
    Assume first that $n$ is squarefree and it as
    \[
        n=p_1p_2\cdots p_k,\quad p_i\text{ distinct primes}.
    \]
    Let $d$ be a divisor of $n$. For simplicitt of exposition and WLOG we will write $d$ as
    \[
        d=p_1p_2\cdots p_\ell,\quad\text{ and }\quad \frac{n}{d}=o_{\ell+1}p_{\ell+2}\cdots p_{k}.
    \]
    Observe that $\mu(n)=\mu(d)\mu\left(\frac{n}{d}\right)$, since $\mu$ is multiplicative and $\left(d,\frac{n}{d}\right)=1$. Also 
    \[
        \log n=\sum_{d\mid n}\Lambda\left(\frac{n}{d}\right).
    \]
    It follows that
    \begin{align*}
        \mu(n)\log n&=\sum_{d\mid n}\mu(n)\Lambda\left(\frac{n}{d}\right)\\
        &=\sum_{d\mid n}\mu(d)\mu\left(\frac{n}{d}\right)\Lambda\left(\frac{n}{d}\right).   
    \end{align*}
    However, one can verify that $\mu(m)\Lambda(m)=-\Lambda(m)$ for all squarefree $m$. So 
    \begin{align*}
        \mu(n)\log n&=-\sum_{d\mid n}\mu(d)\Lambda\left(\frac{n}{d}\right)\\
        &=-\mu\ast\Lambda(n).
    \end{align*}
    Now if $n$ is not squarefree, then let's assume first that $n$ is divisible by the squares of at least two primes. In this case
    \[
        \mu(n)\log n=0
    \]
    and
    \[
        -\sum_{d\mid n}\mu(d)\Lambda\left(\frac{n}{d}\right)=\sum_{p^m\mid n}\mu\left(\frac{n}{p^m}\right)\Lambda(p^m).
    \]
    In fact, if we write 
    \[
        n=p_1^{e_1}p_2^{e_2}\cdots p_k^{e_k}
    \]
    with $e_1,e_2\geq2$, then $\frac{n}{p_i^{e_i}}$ is not squarefree for all $m_i\leq e_i$. Hence, 
    \[
        \sum_{d\mid n}\mu(d)\Lambda\left(\frac{n}{d}\right)=0=\mu(n)\log n.
    \]
    You can follow a similar arguemnt to show that
    \[
        \mu(n)\log n=0=-\sum_{d\mid n}\mu(d)\Lambda\left(\frac{n}{d}\right),
    \]
    in the remaining case where $n$ is divisible by the square of only one prime number, let's call is $p_0$. We write $n=p_0^{e_0}n_0$, with $n_0$ squarefree and $(n_0,p_0)=1$.
\end{proof}
Let's go back and prove lemma 3.8 which states that
\[
    \left|\sum_{n\leq x}\frac{\mu(n)}{n}\right|\leq1.
\]
\begin{proof}[Proof of Lemma 3.8]
    We will assume WLOG that $x\in\N$. Let's set $x=N$. We consider
    \[
        S(N)=\sum_{n\leq N}e(n)
    \]
    with
    \[
        e(n)=
        \begin{cases}
            1&n=1\\
            0&\text{otherwise}.
        \end{cases}
    \]
    One one hand we have $S(N)=1$. On the other hand, we have
    \begin{align*}
        S(N)&=\sum_{n\leq N}e(n)\\
        &=\sum_{n\leq N}\mu\ast1(n)\\
        &=\sum_{n\leq N}\sum_{d\mid n}\mu(d)\\
        &=\sum_{d\leq N}\mu(d)\sum_{m\leq\frac{N}{d}}1\\
        &=\sum_{d\leq N}\mu(d)\floor{\frac{N}{d}}\\
        &=\sum_{d\leq N}\mu(d)\left(\frac{N}{d}-\fract{\frac{N}{d}}\right)\\
        &=N\sum_{d\leq N}\frac{\mu(d)}{d}-\sum_{d\leq N}\mu(d)\fract{\frac{N}{d}}.
    \end{align*} 
    Now we have
    \[
        1=S(N)=N\sum_{d\leq N}\frac{\mu(d)}{d}-\sum_{d\leq N}\mu(d)\fract{\frac{N}{d}}
    \]
    which can be written as 
    \[
        N\sum_{d\leq N}\frac{\mu(d)}{d}=1+\sum_{d\leq N}\mu(d)\fract{\frac{N}{d}}.
    \]
    Observe that
    \begin{align*}
        \sum_{d\leq N}\mu(d)\fract{\frac{N}{d}}&=\sum_{d\leq N-1}\mu(d)\fract{\frac{N}{d}}\\
        &\leq N-1
    \end{align*}
    Hence, 
    \[
        N\left|\sum_{d\leq N}\frac{\mu(d)}{d}\right|\leq 1+N-1=N,
    \]
    and so
    \[
        \left|\sum_{d\leq N}\frac{\mu(d)}{d}\right|\leq1.\qedhere
    \]
\end{proof}
\begin{proof}[Proof of Theorem 3.7]
    We will show that
    \begin{enumerate}[label=(\roman*)]
        \item $\psi(x)\sim x\Rightarrow H(\mu)=0$
        \item $M(\mu)=0\Rightarrow \Psi(x)\sim x.$
    \end{enumerate}
    Let's proceed with (i).\\
    Suppose that $\Psi(x)\sim x$. Then given any $\varepsilon>0$, there exists $x_0>0$ such that
    \[
        \left|\frac{\Psi(x)}{x}-1\right|\leq\frac{\varepsilon}{4}x
    \]
    for all $x\geq x_0$. i.e., 
    \[
        |\Psi(x)-x|\leq\frac{\varepsilon}{4}x
    \]
    for all $x\geq x_0$.\\
    We need to show that $H(\mu)=0$, i.e.,
    \[
        H(x)=\sum_{n\leq x}\mu(n)\log n,
    \]
    then we want to show that
    \[
        \lim_{x\to 0}\frac{H(x)}{x\log x}=0.
    \]
    In other words given $\varepsilon>0$, show that 
    \[
        \left|\frac{H(x)}{x\log x}\right|\leq\varepsilon
    \]
    for all $x$ sufficiently large. We have 
    \begin{align*}
        H(x)&=\sum_{n\leq x}\mu(n)\log n\\
        &=-\sum_{n\leq x}\sum_{d\mid n}\Lambda\left(\frac{n}{d}\right)\mu(d)\\
        &=-\sum_{d\leq x}\mu(d)\sum_{m\leq\frac{n}{d}}\Lambda(m)\\
        &=-\sum_{d\leq x}\mu(d)\Psi\left(\frac{x}{d}\right)\\
        &=-\sum_{d\leq\frac{x}{x_0}}\Psi\left(\frac{x}{d}\right)-\sum_{\frac{x}{x_0}<d\leq x}\mu(d)\Psi\left(\frac{x}{d}\right)
    \end{align*}
    for $d\leq\frac{x}{x_0}$ we have 
    \[
        \Psi\left(\frac{x}{d}\right)=\Psi\left(\frac{x}{d}\right)-\frac{x}{d}+\frac{x}{d}
    \]
    for $\frac{x}{x_0}<d\leq x$, we have 
    \begin{align*}
        \Psi\left(\frac{x}{d}\right)&=\sum_{n\leq\frac{x}{d}}\Lambda(n)\\
        &\leq\sum_{n\leq\frac{x}{d}}\log n\\
        &\leq\frac{x}{d}\log\left(\frac{x}{d}\right).
    \end{align*}
    Going back to $H(x)$, we have
    \begin{align*}
        H(x)&=-\sum_{d\leq\frac{x}{x_0}}\mu(d)\Psi\left(\frac{x}{d}\right)-\sum_{\frac{x}{x_0}<d\leq x}\mu(d)\Psi\left(\frac{x}{d}\right)\\
        &=-\sum_{d\leq\frac{x}{x_0}}\mu(d)\left(\Psi\left(\frac{x}{d}\right)-\frac{x}{d}\right)+\sum_{d\leq\frac{x}{x_0}}\mu(d)\frac{x}{d}-\sum_{\frac{x}{x_0}<d\leq x}\mu(d)\Psi\left(\frac{x}{d}\right)\\
        \text{so }|H(x)|&\leq\left|\sum_{d\leq\frac{x}{x_0}}\mu(d)\left(\Psi\left(\frac{x}{d}\right)-\frac{x}{d}\right)\right|+\left|x\sum_{d\leq\frac{x}{x_0}}\frac{\mu(d)}{d}\right|+\left|\sum_{\frac{x}{x_0}<d\leq x}\mu(d)\Psi\left(\frac{x}{d}\right)\right|\\
        &\leq\sum_{d\leq\frac{x}{x_0}}\left|\Psi\left(\frac{x}{d}\right)-\frac{x}{d}\right|+\sum_{\frac{x}{x}_0<d\leq x}\left|\Psi\left(\frac{x}{d}\right)\right|+x\left|\sum_{d\leq\frac{x}{x_0}}\frac{\mu(d)}{d}\right|\\
        &\leq\frac{\varepsilon}{4}\sum_{d\leq\frac{x}{x_0}}\frac{x}{d}+\sum_{\frac{x}{x_0}<d\leq x}\frac{x}{d}\log\frac{x}{d}+x\\
        &\leq\frac{\varepsilon}{4}x\left(\log\left(\frac{x}{x_0}\right)+\oh(1)\right)+xx_0\log x_0+x\\
        |H(x)|&\leq\frac{\varepsilon}{4}x\log x+\Oh_{\varepsilon}(x)+xx_0\log x_0+x\\
        \left|\frac{H(x)}{\log x}\right|&\leq\frac{\varepsilon}{4}x+\Oh_{\varepsilon}\left(\frac{1}{\log x}\right)+\frac{x_0\log x_0}{\log x}+\frac{1}{\log x}\\
        &\leq\frac{\varepsilon}{4}+\frac{\varepsilon}{4}+\frac{\varepsilon}{4}+\frac{\varepsilon}{4}\\
        &\leq\varepsilon
    \end{align*}
    for $x$ sufficiently large, as desired.\\
    For (ii), suppose that $M(\mu)$=0, i.e.,
    \[
        \frac{1}{x}\sum_{n\leq x}\mu(n)\to0\text{ as }0x\to\infty.
    \]
    Let's show that
    \[
        \sum_{n\leq x}\Lambda(n)=\Psi(x)\sim x.
    \]
    Recall
    \[
        \log n=\sum_{d\mid n}\Lambda\left(\frac{n}{d}\right),\quad (\log=1\ast\Lambda).
    \]
    By M{\"o}bius inversion
    \begin{align*}
        \Lambda(n)&=\sum_{d\mid n}\log\left(\frac{n}{d}\right)\mu(d)\\
        &=\log\ast\mu(n).
    \end{align*}
    Let $f(n)=d(n)-2\gamma$,
    \begin{align*}
        \sum_{n\leq x}f(n)&=\sum_{n\leq x}d(n)-2\gamma\sum_{n\leq x}1\\
        &=x\log x+(2\gamma-1)x+\oh(\sqrt{x})-2\gamma(x+\oh(1))\\
        \sum_{n\leq x}f(n)&=x\log x-x+\Oh(\sqrt{x}).
    \end{align*}
    Observe that
    \[
        \sum_{n\leq x}\log n=x\log x-x+\Oh(\log x)
    \]
    and so if we get $r(n)=\log n-f(n)$, then 
    \[Z
        R(x)=\sum_{n\leq x}r(n)=\Oh(\sqrt{x}).
    \]
    Now let's get back to
    \begin{align*}
        \Lambda(n)&=\log\ast\mu(n)\\
        &=(r+f)\ast\mu(n)\\
        (r+d(n)-2\gamma\times1)\ast\mu(n)\\
        &=r\ast\mu(n)+d\ast\mu(n)-2\gamma(1\ast\mu)(n).
    \end{align*}
    Hence, 
    \[
        \Lambda(n)=r\ast\mu(m)+1(n)-2\gamma e(n),
    \]
    \[
        \sum_{n\leq x}\Lambda(n)=\sum_{n\leq x}r\ast\mu(n)+\sum_{n\leq x}1(n)-2\gamma\sum_{n\leq x}e(n),
    \]
    \[
        \Psi(x)=E(x)+x+\oh(1)-2\gamma.
    \]
    If we estimate $E(x)$ trivially, we get
    \begin{align*}
        E(x)&=\sum_{n\leq x}r\ast\mu(n)\\
        &=\sum_{d\leq x}\mu(d)\sum_{m\leq\frac{n}{d}}r(m)\\
        |E*(x)|&\leq\sum_{d\leq x}|\mu(d)|\left|\sum_{m\leq\frac{n}{d}}r(m)\right|\\
        &\leq\sum_{d\leq x}\left|\sum_{m\leq\frac{n}{d}}r(m)\right|\\
        &\ll\sum_{d\leq x}\left(\sqrt{\frac{x}{d}}\right)\\
        &=\sqrt{x}\sum_{d\leq x}\left(\frac{1}{\sqrt{d}}\right)\\
        &=\Oh(x).
    \end{align*}
    This is not a good estimate, we want $E(x)=\oh(x)$.\\
    To this end we will need to apply the hyperbola method. We have:
    \begin{align*}
        E(x)&=\sum_{n\leq x}\mu\ast r(n)\\
        &=\sum_{d\leq x}\mu(d)\sum_{m\leq\frac{x}{d}}r(m)\\
        &=I_1+I_2-I_3,
    \end{align*}
    where
    \begin{align*}
        I_1&=\sum_{d\leq\frac{x}{y}}\mu(d)\sum_{m\leq\frac{x}{d}}r(m)=\sum_{d\leq\frac{x}{y}}\mu(d)R\left(\frac{x}{d}\right)\\
        I_2&=\sum_{m\leq y}r(m)\sum_{d\leq\frac{x}{m}}\mu(d)=\sum_{m\leq y}r(m)M_{\mu}\left(\frac{x}{m}\right)\\
        I_3&=\sum_{m\leq y}r(m)\sum_{d\leq\frac{x}{y}}\mu(d)=R(y)M_{\mu}\left(\frac{x}{y}\right).
    \end{align*}
    Let $\varepsilon>0$ be given. There exists $x_0>0$ such that
    \[
        \left|\frac{M_{\mu}(x)}{x}\right|\leq\frac{\varepsilon}{k}
    \]
    for all $x\geq x_0$, here $k$ can be any positive integer. We also know that $R(x)=\Oh(x)$ i.e. $|R(x)|\leq C\sqrt{x}$ for some $C>0$. It follows that 
    \begin{align*}
        |I_1|&=\left|\sum_{d\leq\frac{x}{y}}\mu(d)R\left(\frac{x}{d}\right)\right|\\
        &\leq C\sum_{d\leq\frac{x}{y}}\sqrt{\frac{x}{d}}\\
        &=C\sqrt{x}\sum_{d\leq\frac{x}{y}}\frac{1}{\sqrt{d}}\\
        &\leq C_1\frac{x}{\sqrt{y}}.
    \end{align*}
    We have
    \[
        \frac{|I_1|}{x}\leq\frac{C_1}{\sqrt{y}}.
    \]
    For $I_2$, we have
    \begin{align*}
        |I_2|=\left|\sum_{m\leq y}r(m)M\left(\frac{x}{m}\right)\right|\\
        &\leq\sum_{m\leq y}|r(m)|\left|\left(\frac{x}{m}\right)\right|\\
        &\leq\sum_{m\leq y}|r(m)|\frac{\varepsilon}{k}\left(\frac{x}{m}\right)\\
        &=\frac{\varepsilon}{k}x\sum_{m\leq y}\frac{|r(m)|}{m}C(y),
    \end{align*}
    i.e. 
    \[
        \frac{|I_2|}{x}\leq\frac{\varepsilon}{k}C(y).
    \]
    Finally 
    \begin{align*}
        |I_3|&\leq C\sqrt{y}\frac{\varepsilon}{k}\left(\frac{x}{y}\right)\\
        \frac{|I_3|}{x}&\leq\frac{C\varepsilon}{k\sqrt{y}}.
    \end{align*}
    So far, we have
    \begin{align*}
        \left|\frac{E(x)}{x}\right|\leq\left|\frac{I_1}{x}\right|+\left|\frac{I_2}{x}\right|+\left|\frac{I_3}{x}\right|\\
        &\leq\frac{C_1}{\sqrt{y}}+\frac{\varepsilon}{k}C(y)+\frac{C\varepsilon}{k\sqrt{y}}.
    \end{align*}
    Choosing $y$ large enough, we get that
    \[
        \left|\frac{E(x)}{x}\right|\leq\frac{\varepsilon}{3}+\frac{\varepsilon}{3}+\frac{\varepsilon}{3}=\varepsilon
    \]
    for all $x$ sufficiently large.
\end{proof}
\chapter{Dirichlet Series }
\end{document}